\begin{bt}%[2D6H2-2]
% nguon: Archive/2024/0. [LVH] BG - 12 New - 4in1-20240625T224301Z-001/0. [LVH] BG - 12 New - 4in1/2. [LVH] TEX SGK 12 NEW-20240418T074833Z-001/2. [LVH] TEX SGK 12 NEW/[LVH] SGK-12-KNTT/data/C6/12-SGK-KNTT-6-2.tex #1130

	Trong quân sự, một máy bay chiến đấu của đối phương có thể xuất hiện ở vị trí X với xác suất 0{,}55. Nếu máy bay đó không xuất hiện ở vị trí X thì nó xuât hiện ở vị trí Y. Để phòng thủ, các bệ phóng tên lửa được bố trí tại các vị trí X và Y. Khi máy bay đối phương xuất hiện ở vị trí X hoặc Y thì tên lửa sẽ được phóng để hạ máy bay đó.\\
	Xét phương án tác chiến sau: Nếu máy bay xuất hiện tại X thì bắn 2 quả tên lửa và nếu máy bay xuất hiện tại Y thì bắn 1 quả tên lửa.\\
	Biết rằng, xác suất bắn trúng máy bay của mỗi quả tên lửa là 0{,}8 và các bệ phóng tên lửa hoạt động độc lập. Máy bay bị bắn hạ nếu nó trúng ít nhất 1 quả tên lửa. Tính xác suất bắn hạ máy bay đối phương trong phương án tác chiến nêu trên.
	
	\loigiai{Xét biến cố $A$: ``Máy bay xuất hiện ở vị trí X'', điều đó có nghĩa là biến cố $\overline{A}$: ``Máy bay xuất hiện ở vị trí Y''. Xét biến cố $B$: ``Máy bay bị bắn hạ''. \\
	Ta có $P(B)=P(A)\cdot P(B \mid A)+P(\overline{A}) \cdot P(B \mid \overline{A})$.
	\begin{itemize}
	\item Tính $P(A)$, $P(\overline{A})$: $P(A)=0{,}55$ và $P(\overline{A})=0{,}35$.
	\item Tính $P(B\mid A)$: Đây là xác suất để máy bay bị bắn hạ tại vị trí X. Máy bay bị bắn hạ nếu nó trúng ít nhất một 1 quả tên lửa (trong 2 quả tên lửa đối với máy bay ở vị trí X), mà xác suất bắn trúng máy bay của mỗi quả tên lửa là $0{,}8$, vậy $P(B\mid A)=1-\left(1-0{,}8\right)\left(1-0{,}8\right)=0{,}96$.
	\item Tính $P(B\mid \overline{A})$: Đây là xác suất để máy bay bị bắn hạ tại vị trí Y. Máy bay bị bắn hạ nếu nó trúng ít nhất một 1 quả tên lửa (trong 1 quả tên lửa đối với máy bay ở vị trí Y), mà xác suất bắn trúng máy bay của mỗi quả tên lửa là $0{,}8$, vậy $P(B\mid \overline{A})=0{,}8$.
	\end{itemize}
Vậy $P(B)=P(A)\cdot P(B \mid A)+P(\overline{A}) \cdot P(B \mid \overline{A})=0{,}55\cdot 0{,}96+0{,}35\cdot 0{,}8=0{,}808$. \\
Vậy xác suất để máy bay bị bắn hạ là $0{,}808$.
	}
\end{bt}

\begin{bt}%[2D6H2-2]
% nguon: Archive/2024/0. [LVH] BG - 12 New - 4in1-20240625T224301Z-001/0. [LVH] BG - 12 New - 4in1/2. [LVH] TEX SGK 12 NEW-20240418T074833Z-001/2. [LVH] TEX SGK 12 NEW/[LVH] SGK-12-KNTT/data/C6/12-SGK-KNTT-6-2.tex #1131

	Có hai chuồng thỏ. Chuồng I có 5 con thỏ đen và 10 con thỏ trắng. Chuồng II có 7 con thỏ đen và 3 con thỏ trắng. Trước tiên, từ chuồng II lấy ra ngẫu nhiên 1 con thỏ rồi cho vào chuồng I. Sau đó, từ chuồng I lấy ra ngẫu nhiên 1 con thỏ. Tính xác suất để con thỏ được lấy ra là con thỏ trắng.
	\loigiai{Xét biến cố $A$: ``Con thỏ được lấy ra từ chuồng II để cho vào chuồng I là con thỏ trắng''. 	Xét biến cố $B$: ``Con thỏ được lấy ra từ chuồng I là con thỏ trắng''. \\		
	Ta có $P(B)=P(A)\cdot P(B \mid A)+P(\overline{A}) \cdot P(B \mid \overline{A})$.
\begin{itemize}
	\item Tính $P(A)$: Đây là xác suất để lấy ra ngẫu nhiên 1 con thỏ trắng từ chuồng II rồi cho vào chuồng I. Có $n\left(\Omega\right)=\mathrm{C}^1_{10}$, $n\left(A\right)=\mathrm{C}^1_3$. Vậy $P(A)=\dfrac{3}{10}$.
	\item Tính $P(\overline{A})$: $P(\overline{A})=1-P(A)=\dfrac{7}{10}$.
	\item Tính $P(B\mid A)$: Đây là xác suất để lấy ra ngẫu nhiên 1 con thỏ trắng từ chuồng I với điều kiện đã chọn ra 1 con thỏ trắng từ chuồng II rồi cho vào chuồng I, tức là có 5 con thỏ đen và 11 con thỏ trắng ở trong chuồng I. Tương tự như trên ta có $P(B\mid A)=\dfrac{11}{16}$.
	\item Tính $P(B\mid \overline{A})$: Đây là để lấy ra ngẫu nhiên 1 con thỏ trắng từ chuồng I với điều kiện đã chọn ra 1 con thỏ đen từ chuồng II rồi cho vào chuồng I, tức là có 6 con thỏ đen và 10 con thỏ trắng ở trong chuồng I. Tương tự như trên  ta có $P(B\mid \overline{A})=\dfrac{10}{16}$.
\end{itemize}
Vậy $P(B)=P(A)\cdot P(B \mid A)+P(\overline{A}) \cdot P(B \mid \overline{A})=\dfrac{3}{10}\cdot \dfrac{11}{16}+\dfrac{7}{10}\cdot \dfrac{10}{16}=\dfrac{103}{160}=0{,}64375$. \\
Vậy xác suất để con thỏ được lấy ra là con thỏ trắng là $0{,}64375$.
	}
\end{bt}

\begin{bt}%[2D6H2-2]
% nguon: Archive/2024/0. [LVH] BG - 12 New - 4in1-20240625T224301Z-001/0. [LVH] BG - 12 New - 4in1/2. [LVH] TEX SGK 12 NEW-20240418T074833Z-001/2. [LVH] TEX SGK 12 NEW/[LVH] SGK-12-KNTT/data/C6/12-SGK-KNTT-6-2.tex #1132

	Tại nhà máy X sản xuất linh kiện điện tử tỉ lệ sản phẩm đạt tiêu chuẩn là $80 \%$. Trước khi xuất xưởng ra thị trường, các linh kiện điện tử đều phải qua khâu kiểm tra chất lượng để đóng dấu OTK. Vì sự kiểm tra không tuyệt đối hoàn hảo nên nếu một linh kiện điện tử đạt tiêu chuẩn thì nó có xác suất 0{,}99 được đóng dấu OTK; nếu một linh kiện điện tử không đạt tiêu chuẩn thì nó có xác suất 0{,}95 không được đóng dấu OTK. Chọn ngẫu nhiên một linh kiện điện tử của nhà máy X trên thị trường.
	\begin{enumerate}[a)]
		\item Tính xác suất để linh kiện điện tử đó được đóng dấu OTK.
		\item Dùng sơ đồ hình cây, hãy mô tả cách tính xác suất để linh kiện điện tử được chọn không được đóng dấu OTK.
	\end{enumerate}
	
	\loigiai{Gọi $A$ là biến cố: ``Linh kiện điện tử đó đạt tiêu chuẩn''. Gọi $B$ là biến cố: ``Linh kiện điện tử đó được đóng dấu OTK''.
		\begin{enumerate}[a)]
			\item Ta có $P(B)=P(A)\cdot P(B \mid A)+P(\overline{A}) \cdot P(B \mid \overline{A})$.
\begin{itemize}
	\item Tính $P(A)$: Đây là xác suất để linh kiện đó đạt tiêu chuẩn. Vậy $P(A)=0{,}8$.
	\item Tính $P(\overline{A})$: $P(\overline{A})=1-P(A)=0{,}2$.
	\item Tính $P(B\mid A)$: Đây là xác suất để linh kiện điện tử đó được đóng dấu OTK với điều kiện nó đạt tiêu chuẩn. Vậy $P(B\mid A)=0{,}99$.
	\item Tính $P(B\mid \overline{A})$: Đây là xác suất để linh kiện điện tử đó được đóng dấu OTK với điều kiện nó không đạt tiêu chuẩn. Vậy $P(B\mid \overline{A})=1-0{,}95=0{,}05$.
\end{itemize}
Vậy $P(B)=P(A)\cdot P(B \mid A)+P(\overline{A}) \cdot P(B \mid \overline{A})=0{,}8\cdot 0{,}99+0{,}2\cdot 0{,}05=0{,}802$. \\
Vậy xác suất để  linh kiện điện tử đó được đóng dấu OTK là $0{,}802$.
	\item Ta có sơ đồ hình cây
	\begin{center}
			\begin{tikzpicture}[scale=0.8, font=\footnotesize, line join=round, line cap=round, >=stealth]
		\draw (0,0)--(-2,-2.5) node [midway, shift=(140:4mm)] {$0{,}8$} --(-3,-5) node [midway, shift=(160:4mm)] {$0{,}99$};
		\draw (0,0)--(2,-2.5) node [midway, shift=(40:4mm)] {$0{,}2$} --(3,-5) node [midway, shift=(20:4mm)] {$0{,}95$};
		\draw (-2,-2.5)--(-1,-5) node [midway, shift=(20:4mm)] {$0{,}01$};
		\draw (2,-2.5)--(1,-5) node [midway, shift=(160:4mm)] {$0{,}05$};
		\draw[draw=none,fill=yellow] (0,0) circle [radius=7pt] node[shift=(90:6mm)] {Gốc $O$};
		\draw[draw=none,fill=cyan] (-2,-2.5) circle [radius=7pt] node[shift=(140:6mm)] {$A$};
		\draw[draw=none,fill=green] (2,-2.5) circle [radius=7pt] node[shift=(40:6mm)] {$\overline{A}$};
		\draw[draw=none,fill=magenta] (1,-5) circle [radius=7pt] node[shift=(-90:6mm)] {$B$};
		\draw[draw=none,fill=magenta!40!black!30] (-1,-5) circle [radius=7pt] node[shift=(-90:6mm)] {$\overline{B}$};
		\draw[draw=none,fill=magenta] (-3,-5) circle [radius=7pt] node[shift=(-90:6mm)] {$B$};
		\draw[draw=none,fill=magenta!40!black!30] (3,-5) circle [radius=7pt] node[shift=(-90:6mm)] {$\overline{B}$};
	\end{tikzpicture}
	\end{center}
	Có hai nhánh cây đi từ $O$ tới $\overline{B}$ là $OA\overline{B}$ và $O\overline{AB}$. Vậy xác suất để linh kiện điện tử được chọn không được đóng dấu OTK là $$P(\overline{B})=0{,}8\cdot 0{,}01+0{,}2\cdot 0{,}95=0{,}198.$$
		\end{enumerate}
	}
\end{bt}

\begin{bt}%[2D6V2-2]
% nguon: Archive/2024/0. [LVH] BG - 12 New - 4in1-20240625T224301Z-001/0. [LVH] BG - 12 New - 4in1/2. [LVH] TEX SGK 12 NEW-20240418T074833Z-001/2. [LVH] TEX SGK 12 NEW/[LVH] SGK-12-KNTT/data/C6/12-SGK-KNTT-6-2.tex #1133

	Có hai đội thi đấu môn Bắn súng. Đội I có 5 vận động viên, đội II có 7 vận động viên. Xác suất đạt huy chương vàng của mỗi vận động viên đội I và đội II tương ứng là 0{,}65 và 0{,}55. Chọn ngẫu nhiên một vận động viên.
	\begin{enumerate}[a)]
		\item Tính xác suất để vận động viên này đạt huy chương vàng;
		\item Giả sử vận động viên được chọn đạt huy chương vàng. Tính xác suất để vận động viên này thuộc đội I.
	\end{enumerate}
	\loigiai{
		\begin{enumerate}
		\item Xét biến cố $A$: ``Vận động viên này thuộc đội I''. Xét biến cố $B$: ``Vận động viên này đạt huy chương vàng''.
		Ta có $P(B)=P(A)\cdot P(B \mid A)+P(\overline{A}) \cdot P(B \mid \overline{A})$.
	\begin{itemize}
		\item Tính $P(A)$: Đây là xác suất để vận động viên đó thuộc đội I. Vậy $P(A)=\dfrac{5}{12}$.
		\item Tính $P(\overline{A})$: $P(\overline{A})=1-P(A)=\dfrac{7}{12}$.
		\item Tính $P(B\mid A)$: Đây là xác suất để vận động viên thuộc đội I đạt huy chương vàng. Vậy $P(B\mid A)=0{,}65$.
		\item Tính $P(B\mid \overline{A})$: Đây là xác suất để vận động viên thuộc đội I đạt huy chương vàng. Vậy $P(B\mid \overline{A})=0{,}55$.
	\end{itemize}
	Vậy $P(B)=P(A)\cdot P(B \mid A)+P(\overline{A}) \cdot P(B \mid \overline{A})=\dfrac{5}{12}\cdot 0{,}65+\dfrac{7}{12}\cdot 0{,}55=\dfrac{71}{120}\approx 0{,}59$. \\
	Vậy xác suất để vận động viên này đạt huy chương vàng là khoảng $0{,}59$.

	\item Ta cần tính $P(A\mid B)$. Theo công thức Bayes ta có
	$$P(A\mid B) = \dfrac{P(A) \cdot P(B\mid A)}{P(B)}=\dfrac{\dfrac{5}{12}\cdot 0{,}65}{\dfrac{71}{120}}=\dfrac{65}{142}\approx 0{,}46.$$
	\end{enumerate}	}
\end{bt}

\begin{bt}%[2D6V2-3]
% nguon: Archive/2024/0. [LVH] BG - 12 New - 4in1-20240625T224301Z-001/0. [LVH] BG - 12 New - 4in1/2. [LVH] TEX SGK 12 NEW-20240418T074833Z-001/2. [LVH] TEX SGK 12 NEW/[LVH] SGK-12-KNTT/data/C6/12-SGK-KNTT-6-2.tex #1134

	Một bộ lọc được sử dụng để chặn thư rác trong các tài khoản thư điện tử. Tuy nhiên, vì bộ lọc không tuyệt đối hoàn hảo nên một thư rác bị chặn với xác suất 0{,}95 và một thư đúng (không phải là thư rác) bị chặn với xác suất 0{,}01. Thống kê cho thấy tỉ lệ thư rác là $3 \%$.
	\begin{enumerate}[a)]
		\item Chọn ngẫu nhiên một thư bị chặn. Tính xác suất để đó là thư rác.
		\item Chọn ngẫu nhiên một thư không bị chặn. Tính xác suất để đó là thư đúng.
		\item Trong số các thư bị chặn, có bao nhiêu phần trăm là thư đúng? Trong số các thư không bị chặn, có bao nhiêu phần trăm là thư rác?
	\end{enumerate}
	\loigiai{Xét biến cố $A$: ``Thư đó là thư rác''. Xét biến cố $B$: ``Thư đó bị chặn''.
	\begin{enumerate}[a)]
	\item Ta cần tính $P(A\mid B)$.
	Áp dụng công thức Bayes
	$$P(A\mid B) = \dfrac{P(A)\cdot P(B\mid A)}{P(A)\cdot P(B\mid A) + P(\overline{A})\cdot P(B \mid \overline{A})}.$$
	\begin{itemize}
		\item Tính $P(A)$: Đây là xác suất thư đó là thư rác. Vậy $P(A)=0{,}03$.
		\item Tính $P(\overline{A})$: $P(\overline{A})=1-P(A)=0{,}97$.
		\item Tính $P(B \mid A)$: Đây là xác suất thư đó là thư rác bị chặn. Vậy $P(B \mid A)=0{,}95$.
		\item Tính $P(B \mid \overline{A})$: Đây là xác suất thư đó là thư đúng bị chặn. Vậy $P(B \mid \overline{A})=0{,}01$.
	\end{itemize}
	Vậy $P(A\mid B) = \dfrac{P(A)\cdot P(B\mid A)}{P(A)\cdot P(B\mid A) + P(\overline{A})\cdot P(B \mid \overline{A})}=\dfrac{0{,}03\cdot 0{,}95}{0{,}03\cdot 0{,}95+0{,}97\cdot 0{,}01}=\dfrac{285}{382}\approx 0{,}746$.\\
	Vậy xác suất để chọn ngẫu nhiên một thư bị chặn mà thư đó là thư rác là khoảng $0{,}746$.
	\item Ta cần tính $P(\overline{A}\mid \overline{B})$.
	Áp dụng công thức Bayes
	$$P(\overline{A}\mid \overline{B}) = \dfrac{P(\overline{A})\cdot P(\overline{B}\mid \overline{A})}{P(\overline{A})\cdot P(\overline{B}\mid \overline{A}) + P(A)\cdot P(\overline{B} \mid A)}.$$
	\begin{itemize}
		\item Tính $P(\overline{B} \mid A)$: Ta có $(\overline{B} \mid A)=1-P(B \mid A)=0{,}05$.
		\item Tính $P(\overline{B} \mid \overline{A})$: Ta có $P(\overline{B} \mid \overline{A})=1-P(B \mid \overline{A})=0{,}99$.
	\end{itemize}
	Vậy $P(\overline{A}\mid \overline{B}) = \dfrac{P(\overline{A})\cdot P(\overline{B}\mid \overline{A})}{P(\overline{A})\cdot P(\overline{B}\mid \overline{A}) + P(A)\cdot P(\overline{B} \mid A)}=\dfrac{0{,}97\cdot 0{,}99}{0{,}97\cdot 0{,}99+0{,}03\cdot 0{,}05}=\dfrac{3201}{3206}\approx 0{,}998$.\\
	Vậy xác suất để chọn ngẫu nhiên một thư không bị chặn mà thư đó là thư đúng là khoảng $0{,}998$.
	\item Trong số các thư bị chặn, có $74{,}6 \%$ là thư rác, có $25{,}4 \%$ là thư đúng. \\
	Trong số các thư không bị chặn, có $0{,}2 \%$ là thư rác, có $99{,}8 \%$ là thư đúng.
	\end{enumerate}
}

\end{bt}

\begin{bt}%[2D6N2-3]
% nguon: Archive/2024/TaiLieuGiangDay/6. K12 - 4in 1 - XsCoDieuKien/data/CHUONG 6 - XAC SUAT/Bai2-Dang2-CongThucBayes.tex #3193
[CD]%[Lê Vũ Hải, BG12-4in1] 
	Cho hai biến cố $A$, $B$ sao cho $\mathrm{P}(A) = 0{,}6$; $\mathrm{P}(B)=0{,}4$; $\mathrm{P}(A|B) = 0{,}3$. Tính $\mathrm{P}(B|A)$.
	\loigiai{
		Áp dụng công thức Bayes, ta có
		\[\mathrm{P}(B|A)=\dfrac{\mathrm{P}(B)\cdot \mathrm{P}(A|B)}{\mathrm{P}(A)}=\dfrac{0{,}4\cdot 0{,}3}{0{,}6}=0{,}2.\]
	}
\end{bt}

\begin{bt}%[2D6H2-3]
% nguon: Archive/2024/TaiLieuGiangDay/6. K12 - 4in 1 - XsCoDieuKien/data/CHUONG 6 - XAC SUAT/Bai2-Dang2-CongThucBayes.tex #3194
[CKP]%[Lê Vũ Hải, BG12-4in1] 
	Kết quả khảo sát tại một xã cho thấy có $20 \%$ cư dân hút thuốc lá. Tỉ lệ cư dân thường xuyên gặp các vấn đề sức khoẻ về đường hô hấp trong số những người hút thuốc lá và không hút thuốc lá lần lượt là $70\%$, $15\%$. Giả sử ta gặp một cư dân của xã, gọi $A$ là biến cố \lq\lq Người đó có hút thuốc lá\rq\rq\, và $B$ là biến cố \lq\lq Người đó thường xuyên gặp các vấn đề sức khoẻ về đường hô hấp\rq\rq.
	\begin{enumerate}
		\item Vẽ sơ đồ hình cây.
		\item Nếu ta gặp một cư dân của xã thì xác suất người đó thường xuyên gặp các vấn đề sức khoẻ về đường hô hấp là bao nhiêu?
		\item  Nếu ta gặp một cư dân của xã thường xuyên gặp các vấn đề sức khoẻ về đường hô hấp thì xác suất người đó có hút thuốc lá là bao nhiêu?
	\end{enumerate}
	\loigiai{
		\begin{enumerate}
			\item Ta có sơ đồ hình cây sau.
			\begin{center}
				\begin{tikzpicture}[scale=.3,>=stealth,every node/.style={shape=rectangle,draw,rounded corners, color=blue, fill=blue!10}]
					%-------------
					\tikzstyle{block} = [rectangle, draw, fill=blue!10, rounded corners, text centered, text width = 10em, minimum height = 2em]
					%-------------
					\node (c1) {Gặp một cư dân};
					\node (c2) [above right = 2cm of c1]{$A$};
					\node (c3) [ below right= 2cm of c1]{$\overline{A}$};
					\node (c4) [above right = 1cm of c2]{$B$} ;
					\node (c5) [below right = 1cm of c2]{$\overline{B}$};
					\node (c6) [ above right =1cm of c3]{$B$};
					\node (c7) [ below right = 1cm of c3]{$\overline{B}$};
					%--------------
					\draw 
					(17.5,15) node[right] {\text{Kết quả}} 
					(20,11.3) node[right] {$AB$} 
					(20,2.5) node[right] {$A\overline{B}$} 
					(20,-2.5) node[right] {$\overline{A}B$} 
					(20,-11.5) node[right] {$\overline{A} \,\,\overline{B}$};
					%--------------
					\draw 
					(25,15) node[right] {\text{Xác suất}} 
					(25,11.3) node[right] {$0{,}14$} 
					(25,2.5) node[right] {$0{,}06$} 
					(25,-2.5) node[right] {$0{,}12$} 
					(25,-11.5) node[right] {$0{,}68$};
					%------------
					\draw[->] (c1.east) --node[above left]{$0{,}2$} (c2.west);
					\draw[->] (c1.east) --node[below left]{$0{,}8$} (c3.west);
					\draw[->] (c2.east) --node[above left]{$0{,}7$} (c4.west);
					\draw[->] (c2.east) --node[below left]{$0{,}3$} (c5.west);
					\draw[->] (c3.east) --node[above left]{$0{,}15$} (c6.west);
					\draw[->] (c3.east) -- node[below left]{$0{,}85$} (c7.west);
				\end{tikzpicture}
			\end{center}
			\item Ta có $\mathrm{P}(B)=\mathrm{P}(A) \cdot \mathrm{P}(B | A)+\mathrm{P}(\overline{A}) \cdot \mathrm{P}\left(B | \overline{A}\right)=0{,}14+0{,}12=0{,}26$.\\
			Vậy nếu ta gặp một cư dân của xã thì xác suất người đó thường xuyên gặp các vấn đề sức khoẻ về đường hô hấp là $26\%$.
			\item Theo công thức Bayes, ta có $\mathrm{P}\left(A | B\right)=\dfrac{\mathrm{P}(A)\cdot \mathrm{P}\left(B| A\right)}{\mathrm{P}(B)}=\dfrac{0{,}14}{0{,}26} \approx 0{,}54$.\\
			Vậy nếu ta gặp một cư dân của xã thường xuyên gặp các vấn đề sức khoẻ về đường hô hấp thì xác suất người đó có hút thuốc lá là khoảng $54\%$.
		\end{enumerate}
	}
\end{bt}

\begin{bt}%[2D6H2-3]
% nguon: Archive/2024/TaiLieuGiangDay/6. K12 - 4in 1 - XsCoDieuKien/data/CHUONG 6 - XAC SUAT/Bai2-Dang2-CongThucBayes.tex #3195
[CD]%[Lê Vũ Hải, BG12-4in1] 
	Giả sử có một loại bệnh mà tỉ lệ người mắc bệnh là $0{,}1\%$. Giả sử có một loại xét nghiệm, mà ai mắc bệnh khi xét nghiệm cũng có phản ứng dương tính, nhưng tỉ lệ phản ứng dương tính giả là $5\%$ (tức là trong số những người không bị bệnh có $5\%$ số người xét nghiệm lại có phản ứng dương tính).
	\begin{enumerate}
		\item Vẽ sơ đồ cây biểu thị tình huống trên.
		\item Khi một người xét nghiệm có phản ứng dương tính thì khả năng mắc bệnh của người đó là bao nhiêu phần trăm (làm tròn kết quả đến hàng phần trăm).
	\end{enumerate}
	\loigiai{
		\begin{enumerate}
			\item Xét hai biến cố
			\begin{itemize}
				\item[] $K$: \lq\lq Người được chọn ra không mắc bệnh\rq\rq.
				\item[] $D$: \lq\lq Người được chọn ra có phản ứng dương tính\rq\rq. 
			\end{itemize}
			Do tỉ lệ người mắc bệnh là là $0,1\%=0{,}001$ nên $\mathrm{P}(K)=1-0{,}001 = 0{,}999$.\\
			Trong số những người không mắc bệnh có $5\%$ số người có phản ứng dương tính nên \break $\mathrm{P}(D|K) = 5\% = 0{,}05$. Vì ai mắc bệnh khi xét nghiệm cũng có phản ứng dương tính nên $\mathrm{P}\left(\overline{K}\right)=1$.\\
			Ta có sơ đồ cây biểu thị tình huống đã cho
			\begin{center}
				\begin{tikzpicture}[->,>=stealth,line join=round,line cap=round,font=\footnotesize,scale=1]
					\def\xmot{4}
					\def\xhai{8}
					\node (O) at (0,0){};
					\node (B) at (\xmot,1){$K$};
					\node (B1) at (\xmot,-1){$\overline{K}$};
					\node (BA) at (\xhai,2){$D$};
					\node (BA1) at (\xhai,0.3){$\overline{D}$};
					\node (B1A) at (\xhai,-0.3){$D$};
					\node (B1A1) at (\xhai,-1.75){$\overline{D}$};
					
					\foreach \x/\y/\p/\l in 
					{
						O/B/above/$\mathrm{P}(K)=0{,}999$,
						B/BA/above/$\mathrm{P}(D|K)=0{,}05$,
						B/BA1//,
						O/B1/below/$\mathrm{P}\left(\overline{K}\right)=0{,}001$,
						B1/B1A/above/$\mathrm{P}\left(D|\overline{K}\right)=1$,
						B1/B1A1//
					}
					{
						\draw[->] (\x)--(\y)node[midway,\p,scale=0.8,sloped]{\l};
					}
				\end{tikzpicture}
			\end{center}
			\item Ta thấy rằng khả năng mắc bệnh của một người xét nghiệm có phản ứng dương tính chính là $\mathrm{P}\left(\overline{K}|D\right)$. Áp dụng công thức Bayes, ta có:
			\[\mathrm{P}\left(\overline{K}|D\right)=\dfrac{\mathrm{P}\left(\overline{K}\right)\cdot \mathrm{P}\left(D|\overline{K}\right)}{\mathrm{P}\left(\overline{K}\right)\cdot \mathrm{P}\left(D|\overline{K}\right) + \mathrm{P}(K)\cdot\mathrm{P}(D|K)}=\dfrac{0{,}001\cdot 1}{0{,}001\cdot 1 + 0{,}999 \cdot 0{,}05}\approx 1{,}96\%.\]
			Vậy xác suất mắc bệnh của một người xét nghiệm có phản ứng dương tính là $1{,}96\%$.
		\end{enumerate}
	}
\end{bt}

\begin{bt}%[2D6H2-3]
% nguon: Archive/2024/TaiLieuGiangDay/6. K12 - 4in 1 - XsCoDieuKien/data/CHUONG 6 - XAC SUAT/Bai2-Dang2-CongThucBayes.tex #3196
[CTST]%[Lê Vũ Hải, BG12-4in1] 
	Một nhà máy có hai phân xưởng I và II. Phân xưởng I sản xuất $40\%$ số sản phẩm 
	và phân xưởng II sản xuất $60\%$ số sản phẩm. Tỉ lệ sản phẩm bị lỗi của phần xưởng I 
	là $2\%$ và của phân xưởng II là $1\%$.
	\begin{listEX}
		\item  Kiểm tra ngẫu nhiên 1 sản phẩm của nhà máy và tính xác suất để sản phẩm đó bị lỗi. 
		\item Biết rằng sản phẩm được chọn bị lỗi. Hỏi xác suất sản phẩm đó do phân xưởng nào
		sản xuất cao hơn?
	\end{listEX}
	\loigiai{\begin{listEX}
			\item  Gọi $A$ là biến cố “Sản phẩm bị lỗi” và $B$ là biến cố “Sản phẩm lấy ra do phân xưởng I 
			sản xuất”.\\
			Do phân xưởng I sản xuất $40\%$ số sản phẩm và phân xưởng II sản xuất $60\%$ số sản phẩm nên
			$$\mathrm{P}(B)=0{,}4 \text{ và } \mathrm{P}(\overline{B})=1-0{,}4=0{,}6.$$
			Do tỉ lệ sản phẩm bị lỗi của phân xưởng I là $2\%$ và của phân xưởng II là $1\%$ nên
			$$\mathrm{P}(A|B)=0{,}02 \text{ và } \mathrm{P}(A|\overline{B})=0{,}01.$$
			Xác suất để sản phẩm lấy ra bị lỗi là
			$$\mathrm{P}(A)=\mathrm{P}(B)\mathrm{P}(A|B)+\mathrm{P}(\overline{B})\mathrm{P}(A|\overline{B})=0{,}4\cdot0{,}02+0{,}6 \cdot 0{,}01=0{,}014.$$
			\item Nếu sản phẩm lấy ra bị lỗi thì xác suất sản phẩm đó do phân xưởng I sản xuất là 
			$$\mathrm{P}(B|A)=\dfrac{\mathrm{P}(B)\mathrm{P}(A|B)}{\mathrm{P}(A)}=\dfrac{0{,}4 \cdot 0{,}02}{0{,}014}=\dfrac{4}{7}.$$
			Nếu sản phẩm lấy ra bị lỗi thì xác suất sản phẩm đó do phân xưởng II sản xuất là 
			$$\mathrm{P}(\overline{B}|A)=1-\mathrm{P}(B|A)=\dfrac{3}{7}.$$
			Vậy nếu sản phẩm lấy ra bị lỗi thì xác suất sản phẩm đó do phân xưởng I sản xuất cao hơn 
			xác suất sản phẩm đó do phân xưởng II sản xuất.
	\end{listEX}}
\end{bt}

\begin{bt}%[2D6H2-3]
% nguon: Archive/2024/TaiLieuGiangDay/6. K12 - 4in 1 - XsCoDieuKien/data/CHUONG 6 - XAC SUAT/Bai2-Dang2-CongThucBayes.tex #3200
[CKP]%[Lê Vũ Hải, BG12-4in1] 
	Bạn Nam tham gia một gian hàng trò chơi dân gian trong hội xuân của trường. Trò chơi có hai lượt chơi. Xác suất để Nam thắng ở lượt chơi thứ nhất là $0{,}6$. Nếu Nam thắng ở lượt chơi thứ nhất thì xác suất Nam thắng ở lượt chơi thứ hai là $0{,}8 $. Ngược lại, nếu Nam thua ở lượt chơi thứ nhất thì xác suất Nam thắng ở lượt chơi thứ hai là $0{,}3$.
	\begin{listEX}
		\item Vẽ sơ đồ hình cây mô tả các khả năng xảy ra và xác suất tương ứng khi Nam tham gia trò chơi này.
		\item 	Biết Nam đã thắng ở lượt chơi thứ hai, tính xác suất Nam thắng ở lượt chơi thứ nhất.
	\end{listEX}
	\loigiai{
		\begin{listEX}
			\item Sơ đồ hình cây
			\begin{center}
				\begin{tikzpicture}[scale=.3,>=stealth,every node/.style={color=blue}]				
					%-------------
					\tikzstyle{block} = [rectangle, draw, fill=blue!20, rounded corners, text centered, text width = 10em, minimum height = 2em]
					%-------------
					\node (c1) [block] {Trò chơi};
					\node (c2) [block, above right = 3cm of c1]{Nam thắng lượt chơi thứ nhất};
					\node (c3) [block, below right= 3cm of c1]{Nam thua lượt chơi thứ nhất };
					\node (c4) [block,above right = 1.5cm of c2]{Nam thắng lượt chơi thứ hai};
					\node (c5) [block,below right = 1.5cm of c2]{Nam thua lượt chơi thứ hai};
					\node (c6) [block, above right =1.5cm of c3]{Nam thắng lượt chơi thứ hai};
					\node (c7) [block, below right = 1.5cm of c3]{Nam thua lượt chơi thứ hai};
					%--------------
					\draw[->] (c1.east) -- (c2.west);
					\draw[->] (c1.east) -- (c3.west);
					\draw[->] (c2.east) -- (c4.west);
					\draw[->] (c2.east) -- (c5.west);
					\draw[->] (c3.east) -- (c6.west);
					\draw[->] (c3.east) -- (c7.west);
					\draw ($(c1.east)!.5!(c2.west)$) node [below right=-.1]{\color{red}$\mathrm{P}(A) =0{,}6$};
					\draw ($(c1.east)!.5!(c3.west)$) node [above right=-.1]{\color{red}$\mathrm{P}(\overline{A}) =0{,}4$};
					\draw ($(c2.east)!.5!(c4.west)$) node [below right=-.1]{\color{red}$\mathrm{P}(B|A) =0{,}8$};
					\draw ($(c2.east)!.5!(c5.west)$) node [above right=-.1]{\color{red}$\mathrm{P}(\overline{B}|A) =0{,}2$};
					\draw ($(c3.east)!.5!(c6.west)$) node [below right=-.1]{\color{red}$\mathrm{P}(B|\overline{A}) =0{,}3$};
					\draw ($(c3.east)!.5!(c7.west)$) node [above right=-.1]{\color{red}$\mathrm{P}(\overline{B}|\overline{A}) =0{,}7$};
				\end{tikzpicture}
			\end{center}
			\item 
			Công thức Bayes cho sự kiện $A$ và $B$ là $\mathrm{P}(A | B)=\dfrac{\mathrm{P}(B | A) \cdot \mathrm{P}(A)}{\mathrm{P}(B)}	$.\\
			Trong đó
			\begin{itemize}
				\item $A$ là Nam thắng ở lượt chơi thứ nhất.
				\item $B$ là Nam thắng ở lượt chơi thứ hai.
				\item $\mathrm{P}(A)$: Xác suất Nam thắng ở lượt chơi thứ nhất $0{,}6$.
				\item $\mathrm{P}(B | A)$: Xác suất Nam thắng ở lượt chơi thứ hai khi đã thắng ở lượt chơi thứ nhất $0{,}8$.
				\item $\mathrm{P}(B | \overline{A})$: Xác suất Nam thắng ở lượt chơi thứ hai khi đã thua ở lượt chơi thứ nhất $0.3$.
			\end{itemize}
			Công thức Bayes:
			$$
			\begin{aligned}[t]
				\mathrm{P}(A | B)&=\dfrac{\mathrm{P}(B | A) \cdot \mathrm{P}(A)}{\mathrm{P}(B | A) \cdot \mathrm{P}(A)+\mathrm{P}(B | \overline{A}) \cdot \mathrm{P}(\overline{A})} \\
				& =\dfrac{0{,}8 \cdot 0{,}6}{0{,}8 \cdot 0{,}6+0{,}3 \cdot 0{,}4} \\
				& = \dfrac{0{,}48}{0{,}48+0{,}12} \\
				& = \dfrac{0{,}48}{0{,}6} \\
				& = 0{,}8.
			\end{aligned}
			$$
			Vậy, xác suất Nam thắng ở lượt chơi thứ nhất khi đã thắng ở lượt chơi thứ hai là khoảng $0{,}8$ hoặc $80 \%$.
		\end{listEX}
	}
\end{bt}

\begin{bt}%[2D6H2-3]
% nguon: Archive/2024/TaiLieuGiangDay/6. K12 - 4in 1 - XsCoDieuKien/data/CHUONG 6 - XAC SUAT/Bai2-Dang2-CongThucBayes.tex #3201
[CKP]%[Lê Vũ Hải, BG12-4in1]  
	Trong một kì sát hạch lái xe có $65 \%$ thí sinh nam. Biết rằng $80 \%$ thí sinh nam và $70\% $ thí sinh nữ đỗ kì sát hạch này.
	\begin{listEX}
		\item Tính tỉ lệ thí sinh đỗ kì sát hạch này.
		\item Chọn ngẫu nhiên một thí sinh đã đỗ kì sát hạch. Tính xác suất thí sinh đó là nữ.
	\end{listEX}
	\loigiai{
		\begin{listEX}
			\item Xét các biến cố sau
			\begin{itemize}
				\item[] $D$: \lq\lq Thí sinh đỗ kì sát hạch\rq\rq 
				\item[] $M$: \lq\lq Thí sinh là nam giới\rq\rq\, 
				\item[] $\overline{M}$: \lq\lq Thí sinh là nữ giới\rq\rq\, 
				\item[] $D|M$: \lq\lq Thí sinh nam đỗ kì sát hạch\rq\rq\, 
				\item[] $D|\overline{M}$: \lq\lq Thí sinh nữ đỗ kì sát hạch\rq\rq\,.
			\end{itemize}
			Theo đề bài ta có
			\begin{center}
				$\mathrm{P}(M) =0{,}65$;  $\mathrm{P}(D|M) =0{,}8$; $\mathrm{P}(\overline{M}) =0{,}35$; $\mathrm{P}(D|\overline{M}) =0{,}7$.
			\end{center}	
			Theo công thức xác suất toàn phần, ta có
			\begin{align*}
				\mathrm{P}(D) &= \mathrm{P}(M) \cdot \mathrm{P}(D|M) + \mathrm{P}(\overline{M}) \cdot \mathrm{P}(D|\overline{M}).
			\end{align*}	
			Thay vào giá trị đã cho
			\begin{align*}
				\mathrm{P}(D) &= 0{,}65 \cdot 0{,}8 + 0{,}35 \cdot 0{,}7 \\
				&= 0{,}765.
			\end{align*}
			\item Xác suất một thí sinh đỗ là nữ\\
			Để tính xác suất này, ta sử dụng công thức Bayes
			\begin{align*}
				\mathrm{P}(\overline{M}|D) &= \dfrac{\mathrm{P}(D|\overline{M}) \cdot \mathrm{P}(\overline{M})}{\mathrm{P}(D)}.
			\end{align*}
			Thay vào giá trị đã cho và kết quả tính được ở trên
			\begin{align*}
				\mathrm{P}(\overline{M}|D) &= \dfrac{0{,}7 \cdot 0{,}35}{0{,}765} \\
				&\approx 0{,}321.
			\end{align*}
			Vậy, xác suất một thí sinh đỗ kì sát hạch là nữ là khoảng $0{,}321$, hay chấp nhận được là khoảng $32{,}1\%$.
		\end{listEX}
	}
\end{bt}

\begin{bt}%[2D6H2-3]
% nguon: Archive/2024/TaiLieuGiangDay/6. K12 - 4in 1 - XsCoDieuKien/data/CHUONG 6 - XAC SUAT/Bai2-Dang2-CongThucBayes.tex #3202
[CD]%[Lê Vũ Hải, BG12-4in1] 
	Năm $2001$, Cộng đồng châu Âu có làm một đợt kiểm tra rất rộng rãi các con bò để phát hiện những con bị bệnh bò điên. Không có xét nghiệm nào cho kết quả chính xác $100 \%$. Một loại xét nghiệm, mà ở đây ta gọi là xét nghiệm $A$ cho kết quả như sau: khi con bò bị bệnh bò điên thì xác suất để có phản ứng dương tính trong xét nghiệm A là $70 \%$ còn khi con bò không bị bệnh thì xác suất để có phản ứng dương tính trong xét nghiệm $A$ là $10 \%$. Biết rằng tỉ lệ bò bị mắc bệnh bò điên ở Hà Lan là $13$ con trên $1~000~000$ con \textit{(Nguồn: F. M. Dekking et al., Amodern introduction to probability and statistics Understanding why and how, Springer, $2005$)}. Hỏi khi một con bò ở Hà Lan có phản ứng dương tính với xét nghiệm $A$ thì xác suất để nó bị mắc bệnh bò điên là bao nhiêu?
	\loigiai{
		Xét hai biến cố\\			
		$N$: \lq\lq Con bò được chọn bị nhiễm bệnh\rq\rq.\\			
		$D$: \lq\lq Con bò được chọn có phản ứng dương tính\rq\rq.\\	
		Khi đó, ta có\\				
		\[\mathrm{P}(N)=\dfrac{13}{1 000 000}=0{,}000013; \qquad \mathrm{P}(\overline{N})=1-\mathrm{P}(N)=0{,}999987;\]		
		\[\mathrm{P}(D|N)=70\%=0{,}7; \qquad \mathrm{P}(D|\overline{N})=10\%=0{,}1.\]
		Áp dụng công thức Bayes, ta có\\
		$\mathrm{P}(N|D)=\dfrac{\mathrm{P}(D|N) \cdot \mathrm{P}(N)}{\mathrm{P}(N) \cdot \mathrm{P}(D|N)+\mathrm{P}(\overline{N})\mathrm{P}(D|\overline{N})}=\dfrac{0{,}7 \cdot 0{,}000013}{0{,}7 \cdot 0{,}000013+0{,}1 \cdot 0{,}999987}\approx 0{,}009\%$.
	}
\end{bt}

\begin{bt}%[2D6H2-3]
% nguon: Archive/2024/TaiLieuGiangDay/6. K12 - 4in 1 - XsCoDieuKien/data/CHUONG 6 - XAC SUAT/Bai2-Dang2-CongThucBayes.tex #3204
[KNTT]%[Lê Vũ Hải, BG12-4in1] 
	Trong một kho rượu có 30\% là rượu loại I. Chọn ngẫu nhiên một chai rượu đưa cho ông Tùng, một người sành rượu, để nếm thử. Biết rằng, một chai rượu loại I có xác suất 0{,}9 để ông Tùng xác nhận là loại I; một chai rượu không phải loại I có xác suất 0{,}95 để ông Tùng xác nhận đây không phải là loại I. Sau khi nếm, ông Tùng xác nhận đây là rượu loại I. Tính xác suất để chai rượu đúng là rượu loại I.
	\loigiai{Xét biến cố $A$: ``Chai rượu đó là chai rượu loại I''. Xét biến cố $B$: ``Ông Tùng xác nhận chai rượu đó là rượu loại I''.\\
		Ta cần tính $\mathrm{P}(A\mid B)$.
		Áp dụng công thức Bayes
		$$\mathrm{P}(A\mid B) = \dfrac{\mathrm{P}(A)\cdot \mathrm{P}(B\mid A)}{\mathrm{P}(A)\cdot \mathrm{P}(B\mid A) + \mathrm{P}(\overline{A})\cdot \mathrm{P}(B \mid \overline{A})}.$$
		\begin{itemize}
			\item Tính $P\mathrm{P}(A)$: Đây là chai rượu đó là chai rượu loại I. Vậy $\mathrm{P}(A)=0{,}3$.
			\item Tính $\mathrm{P}(\overline{A})$: $\mathrm{P}(\overline{A})=1-\mathrm{P}(A)=0{,}7$.
			\item Tính $\mathrm{P}(B \mid A)$: Đây là xác suất ông Tùng xác nhận đúng một chai rượu loại I là một chai rượu loại I. Vậy $\mathrm{P}(B \mid A)=0{,}9$.
			\item Tính $\mathrm{P}(B \mid \overline{A})$: Đây là xác suất ông Tùng xác nhận sai một chai rượu không phải loại I là một chai rượu loại I. Vậy $\mathrm{P}(B \mid \overline{A})=1-0{,}95=0{,}05$.
		\end{itemize}
		Vậy ta có 
		$$\mathrm{P}(A\mid B) = \dfrac{\mathrm{P}(A)\cdot \mathrm{P}(B\mid A)}{\mathrm{P}(A)\cdot \mathrm{P}(B\mid A) + \mathrm{P}(\overline{A})\cdot \mathrm{P}(B \mid \overline{A})}=\dfrac{0{,}3\cdot 0{,}9}{0{,}3\cdot 0{,}9+0{,}7\cdot 0{,}01}=\dfrac{270}{277}\approx 0{,}975.$$
		Vậy xác suất để chai rượu mà ông Tùng xác nhận là rượu loại I đúng là rượu loại I là khoảng $0{,}975$.}
\end{bt}

\begin{bt}%[2D6H2-3]
% nguon: Archive/2024/TaiLieuGiangDay/6. K12 - 4in 1 - XsCoDieuKien/data/CHUONG 6 - XAC SUAT/Bai2-Dang2-CongThucBayes.tex #3205
[CD]%[Lê Vũ Hải, BG12-4in1] 
	Một loại linh kiện do hai nhà máy số I, số II cùng sản xuất. Tỉ lệ phế phẩm của các nhà máy I, II lần lượt là $4\%$; $3\%$. Trong một lô linh kiện để lẫn lộn $80$ sản phẩm của nhà máy số I và $120$ sản phẩm của nhà máy số II. Một khách hàng lấy ngẫu nhiên một linh liện từ lô hàng đó.
	\begin{enumerate}
		\item Tính xác suất để linh kiện được lấy ra là linh kiện tốt.
		\item Giả sử linh kiện được lấy ra là linh kiện phế phẩm. Xác suất linh kiện đó do nhà máy nào sản xuất là cao nhất?
	\end{enumerate}
	\loigiai{
		\begin{enumerate}
			\item Xét các biến cố
			\begin{itemize}
				\item[] $A$: \lq\lq Linh kiện lấy ra là linh kiện tốt\rq\rq.
				\item[] $B_1$: \lq\lq Linh kiện lấy ra là linh kiện từ nhà máy số I\rq\rq.
				\item[] $B_2$: \lq\lq Linh kiện lấy ra là linh kiện từ nhà máy số II\rq\rq.
			\end{itemize}
			Theo đề bài, ta có
			\begin{listEX}[2]
				\item[] $\mathrm{P}(A|B_1) =1 - 0{,}04 = 0{,}96$.
				\item[] $\mathrm{P}\left(A|B_2\right) = 1-0{,}03 = 0{,}97$.
				\item[] $\mathrm{P}(B_1)=\dfrac{80}{200}=0{,}4$.
				\item[] $\mathrm{P}\left(B_2\right) = \dfrac{120}{200}=0{,}6$.
			\end{listEX}
			Khi đó áp dụng công thức xác suất toàn phần, ta có
			\[\mathrm{P}(A) = \mathrm{P}(A|B_1)\cdot \mathrm{P}(B_1) + \mathrm{P}\left(A|B_2\right)\cdot\mathrm{P}\left(B_2\right)=0{,}96\cdot 0{,}4 + 0{,}97\cdot 0{,}6=0{,}966.\]
			\item Ta có $\mathrm{P}\left(\overline{A}\right)=1-\mathrm{P}(A) = 0{,}034$.\\
			Áp dụng công thức Bayes, ta có
			\begin{itemize}
				\item $\mathrm{P}\left(B_1|\overline{A}\right)=\dfrac{\mathrm{P}\left(\overline{A}|B_1\right)\cdot \mathrm{P}(B_1)}{\mathrm{P}\left(\overline{A}\right)} = \dfrac{0{,}04\cdot 0{,}4}{0{,}034} = \dfrac{8}{17}\approx 0{,}048$.
				\item $\mathrm{P}\left(B_2|\overline{A}\right)=\dfrac{\mathrm{P}\left(\overline{A}|B_2\right)\cdot \mathrm{P}(B_2)}{\mathrm{P}\left(\overline{A}\right)} = \dfrac{0{,}03\cdot 0{,}6}{0{,}034} = \dfrac{8}{167}\approx 0{,}054$.
			\end{itemize}
			Vậy với điều kiện linh kiện lấy ra là linh kiện phế phẩm thì xác suất linh kiện đó do nhà máy I sản xuất là cao nhất.
		\end{enumerate}
	}
\end{bt}

\begin{bt}%[2D6V2-3]
% nguon: Archive/2024/TaiLieuGiangDay/6. K12 - 4in 1 - XsCoDieuKien/data/CHUONG 6 - XAC SUAT/Bai2-Dang2-CongThucBayes.tex #3206
[CKP, Mức độ 3]%[Lê Vũ Hải, BG12-4in1] 
	Một nghiên cứu đã chỉ ra rằng tỉ lệ người bị lao phổi trong nhóm $X$ những người mắc phải hội chứng suy giảm miễn dịch $ {H}$ là $15{,}2 \%$. Kết quả nghiên cứu về một số triệu chứng lâm sàng như có ho trong vòng bốn tuần, hoặc có bị sốt trong vòng bốn tuần, hoặc ra mồ hôi ban đêm từ ba tuần trở lên của nhóm $X$ cho thấy.
	\begin{itemize}
		\item Trong số những người mắc bệnh lao phổi, có $93{,}2 \%$ trường hợp có ít nhất một triệu chứng;
		\item Trong số những người không mắc bệnh lao phổi, có $35{,}8 \%$ trường hợp không có triệu chứng nào.
	\end{itemize}
	Nếu bác sĩ gặp một bệnh nhân thuộc nhóm $X$ và bệnh nhân đó có ít nhất một triệu chứng trên thì xác suất bệnh nhân này mắc bệnh lao phổi là bao nhiêu?
	\loigiai{
		Gọi $A$ là biến cố \lq\lq Người bị lao phổi\lq\lq.\\
		$\overline{A}$ là biến cố \lq\lq Người không mắc lao phổi\lq\lq.\\
		$B$ là biến cố \lq\lq Những người có ít nhất một triệu chứng\rq\rq.\\
		$\overline{B}$ là biến cố \lq\lq Những người không có triệu chứng\rq\rq.\\
		Ta có $\mathrm{P}(A)=0{,}152$. \\
		Khi đó, xác suất những người không mắc lao phổi là $$\mathrm{P}\left(\overline{A}\right)=1-0{,}152=0{,}848.$$
		Ta có xác suất những người có ít một triệu chứng trong những người mắc lao phổi là $$\mathrm{P}\left(B|A\right)=0{,}932.$$
		Khi đó, xác suất những người không có triệu chứng trong những người mắc lao phổi là  $$\mathrm{P}\left(\overline{B}|A\right)=1-0{,}932=0{,}068.$$
		Mặt khác, ta có xác suất những người không có triệu chứng  $$\mathrm{P}\left(\overline{B}|\overline{A}\right)=0{,}358.$$
		Khi đó, xác suất những người có ít nhất một triệu chứng trong những người không mắc bệnh lao phổi là $$\mathrm{P}\left(B|\overline{A}\right)=1-0{,}358=0{,}642.$$
		Ta có $\mathrm{P}(B)=\mathrm{P}\left(B|A\right)\cdot \mathrm{P}(A)+\mathrm{P}\left(B|\overline{A}\right)\cdot \mathrm{P}\left(\overline{A}\right)=0{,}932\cdot 0{,}152+0{,}642\cdot 0{,}848=0{,}68608$.\\
		Theo công thức Bayes, ta có $\mathrm{P}(A|B)=\dfrac{\mathrm{P}(A)\cdot\mathrm{P}\left(B|A\right)}{\mathrm{P}(B)}=\dfrac{0{,}152\cdot 0{,}932}{0{,}68608}\approx 0{,}2065$.
	}
\end{bt}

\begin{bt}%[2D6V2-3]
% nguon: Archive/2024/TaiLieuGiangDay/6. K12 - 4in 1 - XsCoDieuKien/data/CHUONG 6 - XAC SUAT/Bai2-Dang2-CongThucBayes.tex #3207
[CKP, Mức độ 3]%[Lê Vũ Hải, BG12-4in1] 
	Có hai hộp đựng các viên bi cùng kích thước và khối lượng. Hộp thứ nhất chứa $5$ viên bi đỏ và $5$ viên bi xanh, hộp thứ hai chứa $6$ viên bi đỏ và $4$ viên bi xanh. Lấy ngẫu nhiên một viên bi từ hộp thứ nhất chuyển sang hộp thứ hai, sau đó lấy ra ngẫu nhiên một viên bi từ hộp thứ hai (giả sử viên bi được lấy ra từ hộp thứ hai là bi đỏ). Tính xác suất viên bi đỏ đó là của hộp thứ nhất.
	\loigiai{
		Gọi $A$ là biến cố \lq\lq Viên bi được lấy ra từ hộp thứ hai là bi đỏ\rq\rq.\\
		$B$ là biến cố \lq\lq Viên bi được lấy ra từ hộp thứ nhất chuyển sang hộp thứ hai là viên bi đỏ\rq\rq.\\
		$\overline{B}$ là biến cố \lq\lq Viên bi được lấy ra từ hộp thứ nhất chuyển sang hộp thứ hai là bi xanh\rq\rq.\\
		Ta có $\mathrm{P}(B)=\dfrac{5}{10}=\dfrac{1}{2}$; $\mathrm{P}(\overline{B})=\dfrac{5}{10}=\dfrac{1}{2}$.\\
		Nếu viên bi được lấy ra từ hộp thứ nhất chuyển sang hộp thứ hai là bi đỏ thì sau khi chuyển, hộp thứ hai có $7$ bi đỏ và $4$ bi xanh. Do đó $\mathrm{P}(A\mid B) =\dfrac{7}{11}$.\\
		Nếu viên bi được lấy ra từ hộp thứ nhất chuyển sang hộp thứ hai là bi xanh thì sau khi chuyển, hộp thứ hai có $6$ bi đỏ và $5$ bi xanh. Do đó $\mathrm{P}(A\mid B)=\dfrac{6}{11}$.\\
		Áp dụng công thức xác suất toàn phần ta có
		$$\mathrm{P})(A) = \mathrm{P}(B) \cdot \mathrm{P}(A\mid B) + \mathrm{P}(\overline{B}) \cdot \mathrm{P}(A\mid \overline{B}) = \dfrac{1}{2} \cdot \dfrac{7}{11} + \dfrac{1}{2}\cdot \dfrac{6}{11} = \dfrac{13}{22}.$$
		Xác suất viên bi được lấy ra từ hộp thứ nhất chuyển sang hộp thứ hai là viên bi đỏ khi biết viên bi lấy ra từ hộp thứ hai là bi đỏ là
		$$\mathrm{P}(B\mid A) =\dfrac{\mathrm{P}(B) \cdot \mathrm{P}(A\mid B)}{\mathrm{P}(A)} = \dfrac{\dfrac{1}{2} \cdot \dfrac{6}{11}}{\dfrac{13}{22}} = \dfrac{6}{13}.$$
	}
\end{bt}

\begin{bt}%[2D6V2-3]
% nguon: Archive/2024/TaiLieuGiangDay/6. K12 - 4in 1 - XsCoDieuKien/data/CHUONG 6 - XAC SUAT/Bai2-Dang2-CongThucBayes.tex #3208
[CD, Mức độ 3]%[Lê Vũ Hải, BG12-4in1] 
	Có hai chiếc hộp, hộp I có $5$ viên bi màu trắng và $5$ viên bi màu đen; hộp II có $6$ viên bi màu trắng và $4$ viên bi màu đen. Các viên bi có cùng kích thước và khối lượng. Lấy ngẫu nhiên đồng thời hai viên bi từ hộp I bỏ sang hộp II. Sau đó lấy ngẫu nhiên một viên bi từ hộp II.
	\begin{enumerate}
		\item Tính xác suất để viên bi được lấy ra là viên bi màu trắng.
		\item Giả sử viên bi được lấy ra là viên bi màu trắng. Tính xác suất viên bi màu trắng đó thuộc hộp I.
	\end{enumerate}
	\loigiai{
		\begin{enumerate}
			\item Xét các biến cố
			\begin{itemize}
				\item[] $A$: \lq\lq Viên bi lấy ra là viên màu trắng\rq\rq.
				\item[] $B_1$: \lq\lq 2 viên bi lấy ra từ hộp I có màu trắng\rq\rq. 
				\item[] $B_2$: \lq\lq 2 viên bi lấy ra từ hộp I có màu đen\rq\rq.
				\item[] $B_3$: \lq\lq 2 viên bi lấy ra từ hộp I có cả hai màu đen trắng\rq\rq.
			\end{itemize}
			Ta có 
			\[\mathrm{P}(B_1) = \dfrac{\mathrm{C}^2_{5}}{\mathrm{C}^2_{10}}=\dfrac{2}{9};\quad 
			\mathrm{P}(B_2) = \dfrac{\mathrm{C}^2_5}{\mathrm{C}^2_{10}}=\dfrac{2}{9};\quad
			\mathrm{P}(B_3) = \dfrac{\mathrm{C}^1_5\cdot\mathrm{C}^1_5}{\mathrm{C}^2_{10}}=\dfrac{5}{9}.
			\]
			Áp dụng công thức xác suất toàn phần, ta có
			\allowdisplaybreaks
			\begin{eqnarray*}
				\mathrm{P}(A) 
				&=& \mathrm{P}(A|B_1)\cdot \mathrm{P}(B_1) + \mathrm{P}(A|B_2)\cdot\mathrm{P}(B_2) + \mathrm{P}(A|B_3)\cdot\mathrm{P}(B_3)\\ 
				&=& \dfrac{\mathrm{C}^2_8}{\mathrm{C}^2_{12}}\cdot \dfrac{2}{9} + \dfrac{\mathrm{C}^2_{6}}{\mathrm{C}^2_{12}}\cdot \dfrac{2}{9} + \dfrac{\mathrm{C}^2_7}{\mathrm{C}^2_{12}}\cdot \dfrac{5}{9}\\
				&=& \dfrac{14}{33}\cdot \dfrac{2}{9} + \dfrac{5}{22}\cdot \dfrac{2}{9} + \dfrac{7}{22}\cdot \dfrac{5}{9}\\
				&=& \dfrac{191}{594}.
			\end{eqnarray*}
			\item Từ yêu cầu, ta cần tìm $\mathrm{P}(B_1|A) + \mathrm{P}(B_3|A)$.\\
			Áp dụng công thức Bayes, ta có
			\allowdisplaybreaks
			\begin{eqnarray*}
				\mathrm{P}(B_1|A) + \mathrm{P}(B_3|A) 
				&=& \dfrac{\mathrm{P}(A|B_1)\cdot \mathrm{P}(B_1)}{\mathrm{P}(A)} + \dfrac{\mathrm{P}(A|B_3)\cdot \mathrm{P}(B_3)}{\mathrm{P}(A)}\\
				&=& \dfrac{\dfrac{14}{33}\cdot \dfrac{2}{9}}{\dfrac{191}{594}} + \dfrac{\dfrac{7}{22}\cdot \dfrac{5}{9}}{\dfrac{191}{594}}\\
				&=& \dfrac{161}{191}.
			\end{eqnarray*}
		\end{enumerate}
	}
\end{bt}

\begin{bt}%[2D6H1-2]
% nguon: Archive/2024/TaiLieuGiangDay/6. K12 - 4in 1 - XsCoDieuKien/data/CHUONG 6 - XAC SUAT/Bai1-Dang2-TinhXacSuatCoDieuKienTheoCongThuc.tex #3219
[CD, Mức độ 1]%[BG-12-New-4in1, Quyết Nguyễn]
	Cho hai biến cố $A$, $B$ có $\mathrm{P}(A)=0,4$; $\mathrm{P}(B)=0,6$; $\mathrm{P}(AB)=0,2$. Tính các xác suất $\mathrm{P}(A\mid B)$ và $\mathrm{P}(B\mid A)$.
	\loigiai{
		Ta có 
		\begin{itemize}
			\item $\mathrm{P}(A\mid B)=\dfrac{\mathrm{P}(AB)}{\mathrm{P}(B)}=\dfrac{0{,}2}{0{,}6}=\dfrac{1}{3}$.
			\item $\mathrm{P}(B|A)=\dfrac{\mathrm{P}(AB)}{\mathrm{P}(A)}=\dfrac{0{,}2}{0{,}4}=0{,}5$.
		\end{itemize}
	}	
\end{bt}

\begin{bt}%[2D6H2-2]
% nguon: Archive/2024/TaiLieuGiangDay/6. K12 - 4in 1 - XsCoDieuKien/data/CHUONG 6 - XAC SUAT/Bai1-Dang2-TinhXacSuatCoDieuKienTheoCongThuc.tex #3223
%[BG-12-New-4in1, Quyết Nguyễn]
	Một phòng nghiên cứu dược học cho $500$ người bị bệnh $H$ dùng hai loại thuốc $X, Y$ để điều trị. Một số người được điều trị bằng thuốc $X$ và số người còn lại được điều trị bằng thuốc $Y$. Kết quả nghiên cứu được trình bày ở bảng $6.2$. 
	\begin{center}
		\begin{tikzpicture}
			\begin{scope}[xscale=4.4]
				\path
				(0,0) foreach \i[count=\k] in {$X$,$Y$} {++(1,0)node(1\k){\i}}
				(0,-1) node {Khỏi bệnh} foreach \i[count=\k] in {$180$,$190$} {++(1,0)node(2\k){\i}}
				(0,-2) node{Không khỏi bệnh} foreach \i[count=\k] in {$60$,$70$} {++(1,0)node(3\k){\i}}
				%(0,-3) node{tiện} foreach \i[count=\k] in {,$O_1$,,$O_4$,} {++(1,0)node(4\k){\i}}
				%(0,-4) node{điện} foreach \i[count=\k] in {,$O_1$,,$O_4$,} {++(1,0)node(5\k){\i}}
				%(0,-5) node{nước} foreach \i[count=\k] in {,$O_1$,,$O_4$,} {++(1,0)node(6\k){\i}}
				;
				%\path
				%(23.south east) node{$\times$}
				%;
				\draw[shift={(-0.5,.5)}] (0,0) grid (3.,-3)
				(0,0)--(1.,-1)
				(0,-1) node[above right]{Tình trạng}
				(1,0) node[below left]{Loại thuốc}
				;
			\end{scope}
		\end{tikzpicture}
	\end{center}
	Chọn ngẫu nhiên một người trong số này. Gọi $A$ là biến cố \lq\lq Người được chọn khỏi bệnh\rq\rq, $B$ là biến cố \lq\lq Người được chọn điều trị bằng thuốc $X$\rq\rq, $C$ là biến cố \lq\lq Người được chọn điều trị bằng thuốc $Y$\rq\rq.
	\begin{enumerate}[a)]
		\item Tính và giải thích ý nghĩa của $\mathrm{P}(A \mid B)$ và $\mathrm{P}(A \mid C)$.
		\item Có thể nói loại thuốc nào có hiệu quả hơn trong việc điều trị bệnh $H$?
	\end{enumerate}
	\loigiai{
		\begin{enumerate}[a)]
			\item $\mathrm{P}(A \mid  B)=\dfrac{n\left(AB\right)}{n(B)}=\dfrac{180}{240}=\dfrac{3}{4}$ và $\mathrm{P}(A \mid C)=\dfrac{n\left(A\cap C\right)}{n(C)}=\dfrac{190}{260}=\dfrac{19}{26}$.  \\
			Theo các kết quả trên, xác suất để một người khỏi bệnh khi được chọn điều trị bằng thuốc $X$ là $\dfrac{3}{4}$ và xác suất để một người khỏi bệnh khi được chọn điều trị bằng thuốc $Y$ là $\dfrac{19}{26}$.
			\item Do $\dfrac{3}{4}>\dfrac{19}{26}$ nên loại thuốc $X$ có hiệu quả hơn loại thuốc $Y$ trong việc điều trị bệnh $H$.
		\end{enumerate}
	}
\end{bt}

\begin{bt}%[2D6H2-2]
% nguon: Archive/2024/TaiLieuGiangDay/6. K12 - 4in 1 - XsCoDieuKien/data/CHUONG 6 - XAC SUAT/Bai1-Dang2-TinhXacSuatCoDieuKienTheoCongThuc.tex #3224
%[BG-12-New-4in1, Quyết Nguyễn]
	Một xí nghiệp dệt may có những dải của một loại vải đang sản xuất theo một quy trình đặc biệt. Những dải này có thể bị lỗi theo hai hướng: lỗi chiều dài và lỗi kết cấu. Thông qua đợt kiểm tra quy trình sản xuất, người ta thấy rằng có $10\%$ dải không đạt yêu cầu về chiều dài, $5\%$ dải không đạt yêu cầu và kết cấu và chỉ có $0{,}8\%$ dải không đạt yêu cầu về cả chiều dài và kết cấu.
	\begin{enumerate}[a)]
		\item Nếu chọn ngẫu nhiên một dải từ quy trình này thì xác suất không đạt yêu cầu về kết cấu là bao nhiêu?
		\item Nếu một dải được chọn ngẫu nhiên từ quy trình này và phép đo nhanh xác định dải đó không đạt yêu cầu về chiều dài, tính xác suất để dải đó không đạt yêu cầu về kết cấu.
	\end{enumerate}
	\loigiai{
		\begin{enumerate}[a)]
			\item Nếu chọn ngẫu nhiên một dải từ quy trình này thì xác suất không đạt yêu cầu về kết cấu là $\dfrac{5}{100}+\dfrac{0{,}8}{100}=\dfrac{29}{500}$.
			\item Gọi $A$ là biến cố \lq\lq Dải được chọn từ quy trình không đạt yêu cầu về kết cấu\rq\rq;\\ 
			$B$ là biến cố \lq\lq Dải được chọn từ quy trình không đạt yêu cầu về chiều dài\rq\rq.\\ 
			Khi đó $AB$ là biến cố \lq\lq Một dải từ quy trình không đạt yêu cầu về cả kết cấu và chiều dài\rq\rq. Ta có $\mathrm{P}(AB)=0{,}8\%=0{,}008$ và $\mathrm{P}(B)=0{,}1$. \\
			Do đó $\mathrm{P}(A\mid B)=\dfrac{\mathrm{P}\left(AB\right)}{\mathrm{P}(B)}=\dfrac{0{,}008}{0{,}1}=0{,}08$.
		\end{enumerate}
	}
\end{bt}

\begin{bt}%[2D6H2-2]
% nguon: Archive/2024/TaiLieuGiangDay/6. K12 - 4in 1 - XsCoDieuKien/data/CHUONG 6 - XAC SUAT/Bai1-Dang2-TinhXacSuatCoDieuKienTheoCongThuc.tex #3225
%[BG-12-New-4in1, Quyết Nguyễn]
	Trong cuộc khảo sát $300$ gia đình ở một khu vực, người ta nhận thấy rằng có $90\%$ gia đình có tivi và $60\%$ gia đình có máy tính bàn. Mỗi gia đình đều có ít nhất một trong hai thiết bị này. Chọn ngẫu nhiên một gia đình. Tính xác suất gia đình có máy tính bàn trong nhóm các gia đình có tivi. 	
	\loigiai{
		Gọi $A$ là biến cố \lq\lq Gia đình được chọn có máy tính bàn\rq\rq; $B$ là biến cố \lq\lq Gia đình được chọn có tivi\rq\rq. Khi đó $AB$ là biến cố \lq\lq Gia đình được chọn có cả máy tính bàn và tivi\rq\rq. \\
		Ta có $n(B)=0{,}9\cdot300=270$ và $n(AB)=0{,}9\cdot300+0{,}6\cdot300-300=150$. \\
		Do đó $\mathrm{P}(A\mid B)=\dfrac{n\left(A\cap B\right)}{n(B)}=\dfrac{150}{270}=\dfrac{5}{9}$.
	}
\end{bt}

\begin{bt}%[2D6H2-2]
% nguon: Archive/2024/TaiLieuGiangDay/6. K12 - 4in 1 - XsCoDieuKien/data/CHUONG 6 - XAC SUAT/Bai1-Dang2-TinhXacSuatCoDieuKienTheoCongThuc.tex #3226
%[BG-12-New-4in1, Quyết Nguyễn]
Trong một lọ có chứa bi đen và bi trắng cùng kích thước và khối lượng, lấy ngẫu nhiên lần lượt hai viên bi ra ngoài và không bỏ vào lại. Biết rằng xác suất để lần đầu lấy được bi đen là $0{,}47$; xác suất để lần đầu lấy được bi đen và lần thứ hai lấy được bi trắng là $0{,34}$. Tính xác suất để lấy được bi trắng ở lần thứ hai với điều kiện lần đầu lấy được bi đen. 
\loigiai{
Gọi $A$ là biến cố \lq\lq Lấy được bi trắng ở lần thứ hai\rq\rq; $B$ là biến cố \lq\lq Lấy được bi đen ở lần đầu\rq\rq. \\
		Do đó $\mathrm{P}(A \mid B)=\dfrac{\mathrm{P}(AB)}{\mathrm{P}(B)}=\dfrac{0{,34}}{0{,47}}=\dfrac{34}{47}$.
	}
\end{bt}

\begin{bt}%[2D6H1-2]
% nguon: Archive/2024/TaiLieuGiangDay/6. K12 - 4in 1 - XsCoDieuKien/data/CHUONG 6 - XAC SUAT/Bai1-Dang1-NhanBietXacSuatCoDieuKien.tex #3237
[CKP, Mức độ 1]%[BG-12-New-4in1, Quyết Nguyễn]
Một hộp có $5$ viên bi cùng kích thước và khối lượng, trong đó có $3$ viên bi màu đỏ và $2$ viên bi màu xanh. Lấy ngẫu nhiên lần lượt $2$ viên bi và không hoàn lại. Tính xác suất để lấy được viên bi thứ hai có màu xanh, biết rằng viên bi thứ nhất có màu đỏ.
\loigiai{
Gọi
\begin{itemize}
\item  $A$ là biến cố \lq\lq Lấy được viên bi thứ hai có màu xanh\rq\rq;
\item  $B$ là biến cố \lq\lq Lấy được viên bi thứ nhất có màu đỏ\rq\rq.
\end{itemize}
Khi đó xác suất để lấy được viên bi thứ hai có màu xanh, biết rằng viên bi thứ nhất có màu đỏ chính là xác suất của $A$ với điều kiện $B$.\\
Vì một viên bi đỏ đã được lấy ra ở lần thứ nhất nên trong hộp còn lại $4$ viên bi, trong đó có $2$ viên bi xanh.\\
Từ đó ta có $\mathrm{P}(A \mid B)=\dfrac{2}{4}=0{,}5$.\\
Vậy xác suất để lấy được viên bi thứ hai có màu xanh, biết rằng viên bi thứ nhất có màu đỏ là $0{,}5$.
}	
\end{bt}

\begin{bt}%[2D6H1-2]
% nguon: Archive/2024/TaiLieuGiangDay/6. K12 - 4in 1 - XsCoDieuKien/data/CHUONG 6 - XAC SUAT/Bai1-Dang1-NhanBietXacSuatCoDieuKien.tex #3243
%[BG-12-New-4in1, Quyết Nguyễn]
	Hộp thứ nhất chứa $4$ viên bi xanh và $3$ viên bi đỏ. Hộp thứ hai chứa $3$ viên bi xanh và $5$ viên bi đỏ. Các viên bi có cùng kích thước và khối lượng. Bạn Thanh lấy ra ngẫu nhiên $1$ viên bi từ hộp thứ nhất bỏ vào hộp thứ hai, sau đó lại lấy ra ngẫu nhiên $1$ viên bi từ hộp thứ hai. Tính xác suất để viên bi lấy ra ở lần thứ hai là viên bi đỏ, biết viên bi lấy ra ở lần thứ nhất là viên bi xanh.
	\loigiai{
		$A$ là biến cố: \lq\lq viên bi lấy ra ở lần thứ hai là viên bi đỏ\rq\rq; $B$ là biến cố: \lq\lq viên bi lấy ra ở lần thứ nhất là viên bi xanh\rq\rq.\\
		Biến cố $B$ xảy ra, nghĩa là lần thứ nhất lấy $1$ viên bi xanh và bỏ vào hộp thứ hai. Khi đó trong hộp thứ hai sẽ có $4$ viên bi xanh và $5$ viên bi đỏ.\\
		Vậy $\mathrm{P}(A\mid B)=\dfrac{5}{9}$. 
	}
\end{bt}

\begin{bt}%[2D6H2-3]
% nguon: Archive/2024/TaiLieuGiangDay/6. K12 - 4in 1 - XsCoDieuKien/data/CHUONG 6 - XAC SUAT/Bai1-Dang34-SoDoCay-CongThucNhanXacSuat.tex #3252
[CD, Mức độ 2]%[BG12-4in1, Võ Thành Huấn]
	Một hộp có $8$ bi màu đỏ và $5$ viên bi màu vàng; các viên bi có kích thước và khối lượng như nhau. Có $5$ viên bi trong hộp được đánh số, trong đó có $3$ viên bi màu đỏ và $2$ viên bi màu vàng. Lấy ngẫu nhiên một viên bi trong hộp. Dùng sơ đồ hình cây, tính xác suất để viên bi được lấy ra, có màu đỏ, biết rằng viên bi đó được đánh số.  
	\loigiai{
		Xét các biến cố sau:
		\begin{itemize}
			\item $A \colon$ \lq\lq Viên bi được lấy ra có đánh số\rq\rq.
			\item $B \colon$ \lq\lq Viên bi được lấy ra có màu đỏ  \rq\rq.
		\end{itemize}
		Khi đó, xác suất để viên bi được lấy ra có màu đỏ, biết ràng viên bi đó được đánh số, chính là xác suất có điều kiện $\mathrm{P}(B\mid A)$.\\
		Sơ đồ hình cây biểu thị cách tính xác suất có điều kiện $\mathrm{P}(B\mid A)$, được vẽ như sau:
		\begin{center}
			\begin{tikzpicture}
				\def\gocm{20}
				\def\gocn{10}
				\def\r{4}
				\tikzset{s/.style={outer sep=0.5 mm,draw=magenta,rectangle,minimum width=2.75cm,rounded corners=1mm}}
				\path(0,0)node(O){}++(\gocm:\r)node[s](A1){A}++(\gocn:\r)node[s](A2){$B$}++(0:\r)node[s](A3){$AB$};
				\path(A1)++({-\gocn}:\r)node[s](a2){$\overline{B}$}++(0:\r)node[s](a3){$A\overline{B}$};
				\path(O)++(-\gocm:\r)node[s](B1){$\overline{A}$}++(\gocn:\r)node[s](B2){$B$}++(0:\r)node[s](B3){$\overline{A}B$};
				\path(B1)++({-\gocn}:\r)node[s](b2){$\overline{B}$}++(0:\r)node[s](b3){$\overline{A}\overline{B}$};
				\foreach \x/\y in {
					O/A1,A1/A2,
					O/B1,B1/B2,
					A1/a2,
					B1/b2}
				\draw[-stealth](\x.east)--(\y.west);
				\path(O)--(A1.west)node[pos=0.5,above,sloped]{$\frac{5}{13}$}(O)--(B1.west)node[pos=0.5,below,sloped]{$\frac{8}{13}$}(B1.east)--(B2.west)node[pos=0.5,above,sloped]{$\frac{5}{8}$}(A1.east)--(A2.west)node[pos=0.5,above,sloped]{$\frac{3}{5}$}
				(A1.east)--(a2.west)node[pos=0.5,below,sloped]{$\frac{2}{5}$}
				(B1.east)--(b2.west)node[pos=0.5,below,sloped]{$\frac{3}{8}$};
				
			\end{tikzpicture}
			
		\end{center}
		Vậy xác suất để viên bi được lấy ra có màu đỏ, biết rằng viên bi đó có đánh số, là $ 0{,}6$.
	}
\end{bt}

\begin{bt}%[2D6H2-3]
% nguon: Archive/2024/TaiLieuGiangDay/6. K12 - 4in 1 - XsCoDieuKien/data/CHUONG 6 - XAC SUAT/Bai1-Dang34-SoDoCay-CongThucNhanXacSuat.tex #3253
[CTST, Mức độ 2]%[BG12-4in1, Võ Thành Huấn]
	Ở một sân bay, người ta sử dụng một loại máy soi tự động phát hiện hàng cấm trong hành lí kí gửi. Máy phát chuông cảnh báo với $95 \%$ các kiện hành lí có chứa hàng cấm và $2 \%$ các kiện hành lí không chứa hàng cấm. Tỉ lệ các kiện hành lí có chứa hàng cấm là $1 \%$.\\
	Chọn ngẫu nhiên một kiện hành lí để soi bằng máy trên. Sử dụng sơ đồ hình cây, tính xác suất của các biến cố:\\
	M: \lq\lq Kiện hành lí có chứa hàng cấm và máy phát chuông cảnh báo \rq\rq;
	N : \lq\lq Kiện hành lí không chứa hàng cấm và máy phát chuông cảnh báo\rq\rq.
	\loigiai
	{
		Gọi $A$ là biến cố \lq\lq Kiện hành lí có chứa hàng cấm\rq\rq và $B$ là biến cố \lq\lq Máy phát chuông cành báo\rq\rq. Ta có
		$$
		\mathrm{P}(B \mid A)=0{,}95 ; \mathrm{P}(B \mid \overline{A})=0{,}02 ; \mathrm{P}(A)=0{,}01.
		$$
		Do đó $\mathrm{P}(\overline{A})=1-\mathrm{P}(A)=0{,}99 ; \mathrm{P}(\overline{B} \mid A)=1-\mathrm{P}(B \mid A)=0{,}05 ; \mathrm{P}(\overline{B} \mid \overline{A})=1-\mathrm{P}(B \mid \overline{A})=0{,}98$.
		Ta có sơ đồ hình cây như sau:
		\begin{center}
			\begin{tikzpicture}
				\def\gocm{20}
				\def\gocn{10}
				\def\r{4}
				\tikzset{s/.style={outer sep=0.5 mm,draw=magenta,rectangle,minimum width=2.75cm,rounded corners=1mm}}
				\path(0,0)node(O){}++(\gocm:\r)node[s](A1){A}++(\gocn:\r)node[s](A2){$B$}++(0:\r)node[s](A3){$AB$}++(0:\r)node[s](A4){$0{,}0095$};
				\path(A1)++({-\gocn}:\r)node[s](a2){$\overline{B}$}++(0:\r)node[s](a3){$A\overline{B}$}++(0:\r)node[s](a4){$0{,}0005$};
				\path(O)++(-\gocm:\r)node[s](B1){$\overline{A}$}++(\gocn:\r)node[s](B2){$B$}++(0:\r)node[s](B3){$\overline{A}B$}++(0:\r)node[s](B4){$0{,}0198$};
				\path(B1)++({-\gocn}:\r)node[s](b2){$\overline{B}$}++(0:\r)node[s](b3){$\overline{A}\overline{B}$}++(0:\r)node[s](b4){$0{,}9702$};
				\foreach \x/\y in {
					O/A1,A1/A2,
					O/B1,B1/B2,
					A1/a2,
					B1/b2}
				\draw[-stealth](\x.east)--(\y.west);
				\path(O)--(A1.west)node[pos=0.5,above,sloped]{$0{,}01$}(O)--(B1.west)node[pos=0.5,below,sloped]{$0{,}99$}(B1.east)--(B2.west)node[pos=0.5,above,sloped]{$0{,}02$}(A1.east)--(A2.west)node[pos=0.5,above,sloped]{$0{,}95$}
				(A1.east)--(a2.west)node[pos=0.5,below,sloped]{$0{,}05$}
				(B1.east)--(b2.west)node[pos=0.5,below,sloped]{$0{,}98$};				
			\end{tikzpicture}			
		\end{center}
		Do $M=A B$ nên $\mathrm{P}(M)=\mathrm{P}(A B)=0{,}0095$.\\
		Do $N=\overline{A} B$ nên $\mathrm{P}(N)=\mathrm{P}(\overline{A} B)=0{,}0198$.
	}
\end{bt}

\begin{bt}%[2D6H2-3]
% nguon: Archive/2024/TaiLieuGiangDay/6. K12 - 4in 1 - XsCoDieuKien/data/CHUONG 6 - XAC SUAT/Bai1-Dang34-SoDoCay-CongThucNhanXacSuat.tex #3254
[CKP]%[BG12-4in1, Võ Thành Huấn]
	Theo kết quả từ trạm nghiên cứu khí hậu tại địa phương $T$, xác suất để một ngày có gió là $0{,}6$; nếu ngày có gió thì xác suất có mưa là $0{,}4$; nếu ngày không có gió thì xác suất có mưa là $0{,}2$. Gọi $G$ là biến cố \lq\lq Ngày có gió\rq\rq~ và $M$ là biến cố \la\la Ngày có mưa\rq\rq.
	\begin{enumerate}
		\item  Vẽ lại sơ đồ hình cây sau và điền vào ô? các giá trị xác suất tương ứng.
		\begin{center}
			\begin{tikzpicture}[node distance=1.5cm, every node/.style={fill=white}, align=center]
				\definecolor{diagram_bg_green}{HTML}{d5e8d4}
				\definecolor{diagram_bg_blue}{HTML}{dae8fc}
				\definecolor{diagram_bg_pink}{HTML}{f8cecc}
				\definecolor{diagram_bd_green}{HTML}{82b366}
				\definecolor{diagram_bd_blue}{HTML}{7494c2}
				\definecolor{diagram_bd_pink}{HTML}{b85450}
				\tikzset{>={Latex[width=2mm,length=2mm]}, base/.style = {rectangle, rounded corners, draw=black, minimum width=1cm, minimum height=1cm, text centered, drop shadow={shadow xshift=0.6mm, shadow yshift=-0.6mm}},
					Style1/.style = {base, fill=diagram_bg_blue, draw=diagram_bd_blue},
					Style2/.style = {base, fill=diagram_bg_pink, draw=diagram_bd_pink},
					Style3/.style = {base, fill=diagram_bg_green, draw=diagram_bd_green},
					Style4/.style = {base, minimum width=2.5cm, fill=orange!15, draw=orange},
				}
				\node (B0) [Style1, text width=1cm] {Ngày};
				\node (B1) [Style2, right of=B0, xshift=2cm, yshift=2cm, text width=1.5cm] {Có gió \\ $(G)$};
				\node (B3) [Style2, right of=B0, xshift=2cm, yshift=-2cm, text width=1.5cm] {Không có gió\\$(\overline{G})$};
				\node (B11) [Style3, right of=B1, xshift=6cm, yshift=1.2cm, text width=7cm] {Có mưa};
				\node (B13) [Style3, right of=B1, xshift=6cm, yshift=-1.2cm, text width=7cm] {Không có mưa};
				\node (B31) [Style3, right of=B3, xshift=6cm, yshift=1.2cm, text width=7cm] {Có mưa};
				\node (B33) [Style3, right of=B3, xshift=6cm, yshift=-1.2cm, text width=7cm] {Không có mưa};
				\draw[->] (B0) -- (B1.west)node[sloped,above,pos=0.5]{$\mathrm{P}(G)=?$};
				\draw[->] (B0) -- (B3.west)node[sloped,below,pos=0.5]{$\mathrm{P}(\overline{G})=?$};
				\draw[->] (B1) -- (B11.west)node[sloped,above,pos=0.5]{$\mathrm{P}(M\mid G)=?$};
				\draw[->] (B1) -- (B13.west)node[sloped,below,pos=0.5]{$\mathrm{P}(\overline{M}\mid G)=?$};
				\draw[->] (B3) -- (B31.west)node[sloped,above,pos=0.5]{$\mathrm{P}(M\mid\overline{G})=?$};
				\draw[->] (B3) -- (B33.west)node[sloped,below,pos=0.5]{$\mathrm{P}(\overline{M}\mid\overline{G})=?$};
			\end{tikzpicture}
		\end{center}
		\item  Tính xác suất $\mathrm{P}(G M)$ và $\mathrm{P}(G \overline{M})$. Nêu ý nghĩa của các xác suất này.
	\end{enumerate}
	\loigiai{
		\begin{enumerate}
			\item Theo đề bài, nếu ngày có gió thì xác suất có mưa là 0{,}4 nên $\mathrm{P}(M \mid G)=0{,}4$.\\ 
			Suy ra $\mathrm{P}(\overline{M} \mid G)=1-0{,}4=0{,}6$.
			Ngày không có gió thì xác suất có mưa là $0{,}2$ nên $\mathrm{P}(M \mid \overline{G})=0{,}2$.\\
			Suy ra $\mathrm{P}(\overline{M} \mid \overline{G})=1-0{,}2=0{,}8$.
			\begin{center}
				\begin{tikzpicture}[node distance=1.5cm, every node/.style={fill=white}, align=center]
					\definecolor{diagram_bg_green}{HTML}{d5e8d4}
					\definecolor{diagram_bg_blue}{HTML}{dae8fc}
					\definecolor{diagram_bg_pink}{HTML}{f8cecc}
					\definecolor{diagram_bd_green}{HTML}{82b366}
					\definecolor{diagram_bd_blue}{HTML}{7494c2}
					\definecolor{diagram_bd_pink}{HTML}{b85450}
					\tikzset{>={Latex[width=2mm,length=2mm]}, base/.style = {rectangle, rounded corners, draw=black, minimum width=1cm, minimum height=1cm, text centered, drop shadow={shadow xshift=0.6mm, shadow yshift=-0.6mm}},
						Style1/.style = {base, fill=diagram_bg_blue, draw=diagram_bd_blue},
						Style2/.style = {base, fill=diagram_bg_pink, draw=diagram_bd_pink},
						Style3/.style = {base, fill=diagram_bg_green, draw=diagram_bd_green},
						Style4/.style = {base, minimum width=2.5cm, fill=orange!15, draw=orange},
					}
					\node (B0) [Style1, text width=1cm] {Ngày};
					\node (B1) [Style2, right of=B0, xshift=2cm, yshift=2cm, text width=1.5cm] {Có gió \\ $(G)$};
					\node (B3) [Style2, right of=B0, xshift=2cm, yshift=-2cm, text width=1.5cm] {Không có gió\\$(\overline{G})$};
					\node (B11) [Style3, right of=B1, xshift=8cm, yshift=1.2cm, text width=7cm] {Có mưa};
					\node (B13) [Style3, right of=B1, xshift=8cm, yshift=-1.2cm, text width=7cm] {Không có mưa};
					\node (B31) [Style3, right of=B3, xshift=8cm, yshift=1.2cm, text width=7cm] {Có mưa};
					\node (B33) [Style3, right of=B3, xshift=8cm, yshift=-1.2cm, text width=7cm] {Không có mưa};
					\draw[->] (B0) -- (B1.west)node[sloped,above,pos=0.5]{$\mathrm{P}(G)=0{,}6$};
					\draw[->] (B0) -- (B3.west)node[sloped,below,pos=0.5]{$\mathrm{P}(\overline{G})=0{,}4$};
					\draw[->] (B1) -- (B11.west)node[sloped,above,pos=0.5]{$\mathrm{P}(M\mid G)=0{,}4$};
					\draw[->] (B1) -- (B13.west)node[sloped,below,pos=0.5]{$\mathrm{P}(\overline{M}\mid G)=0{,}6$};
					\draw[->] (B3) -- (B31.west)node[sloped,above,pos=0.5]{$\mathrm{P}(M\mid\overline{G})=0{,}2$};
					\draw[->] (B3) -- (B33.west)node[sloped,below,pos=0.5]{$\mathrm{P}(\overline{M}\mid\overline{G})=0{,}8$};
				\end{tikzpicture}
			\end{center}
			\item $\mathrm{P}(G M)=\mathrm{P}(G) \cdot \mathrm{P}(M \mid G)=0{,}6 \cdot 0{,}4=0{,}24 $.\\
			$\mathrm{P}(G \overline{M})=\mathrm{P}(G) \cdot \mathrm{P}(\overline{M} \mid G)=0{,}6 \cdot 0{,}6=0{,}36$.\\
			Điểu này có nghĩa là tại địa phương $T$, trong một ngày, xác suất để trời vừa có gió và vừa có mưa là $0{,}24$; xác suất để trời có gió nhưng không có mưa là $0{,}36$.
		\end{enumerate}
	}
\end{bt}

\begin{bt}%[2D6H2-3]
% nguon: Archive/2024/TaiLieuGiangDay/6. K12 - 4in 1 - XsCoDieuKien/data/CHUONG 6 - XAC SUAT/Bai1-Dang34-SoDoCay-CongThucNhanXacSuat.tex #3255
%[BG12-4in1, Võ Thành Huấn]
	Bạn Việt chuẩn bị đi tham quan một hòn đảo trong hai ngày thứ Bảy và Chủ nhật. Ở hòn đảo đó, mỗi ngày chỉ có nắng hoặc mưa, nếu một ngày là nắng thì khả năng xảy ra mưa ở ngày tiếp theo là $20 \%$, còn nếu một ngày là mưa thì khả năng ngày hôm sau vẫn mưa là $30 \%$. Theo dự báo thời tiết, xác suất trời sẽ nắng vào thứ Bảy là $0{,}7$.
	Hãy tìm các giá trị thích hợp thay vào \mbox{?} ở sơ đồ hình cây sau:
	\begin{center}
		\begin{tikzpicture}
			\def\gocm{20}
			\def\gocn{10}
			\def\r{4}
			\tikzset{s/.style={outer sep=0.5 mm,draw=magenta,rectangle,minimum width=2.75cm,rounded corners=1mm}}
			\path(0,0)node(O){}++(\gocm:\r)node[s](A1){Nắng}++(\gocn:\r)node[s](A2){Nắng};
			\path(A1)++({-\gocn}:\r)node[s](a2){Mưa};
			\path(O)++(-\gocm:\r)node[s](B1){Mưa}++(\gocn:\r)node[s](B2){Nắng};
			\path(B1)++({-\gocn}:\r)node[s](b2){Mưa};
			\foreach \x/\y in {
				O/A1,A1/A2,
				O/B1,B1/B2,
				A1/a2,
				B1/b2}
			\draw[-stealth](\x.east)--(\y.west);
			\path(O)--(A1.west)node[pos=0.5,above,sloped]{$\mbox{0{,}7}$}(O)--(B1.west)node[pos=0.5,below]{$\mbox{?}$}(B1.east)--(B2.west)node[pos=0.5,above]{$\mbox{?}$}(A1.east)--(A2.west)node[pos=0.5,above]{$\mbox{?}$}
			(A1.east)--(a2.west)node[pos=0.5,below,sloped]{$\mbox{0{,}2}$}
			(B1.east)--(b2.west)node[pos=0.5,below,sloped]{$\mbox{0{,}3}$};
			%%Node dòng trên
			\path(A2)++(0,1)node{\textbf{Chủ nhật}}++(180:4)node{\textbf{Thứ bảy}};
		\end{tikzpicture}
	\end{center}
\loigiai{
	Gọi $A$ là biến cố \lq\lq Ngày thứ Bảy trời nắng\rq\rq  và $B$ là biến cố \lq\lq Ngày Chủ nhật trời mưa\rq\rq.\\
	Ta có $\mathrm{P}(A)=0{,}7 ; \mathrm{P}(B \mid A)=0{,}2 ; \mathrm{P}(B \mid \overline{A})=0{,}3$.\\
	Do đó $\mathrm{P}(\overline{A})=1-\mathrm{P}(A)=0{,}3 ;\, \mathrm{P}(\overline{B} \mid A)=1-\mathrm{P}(B \mid A)=0{,}8 ;\, \mathrm{P}(\overline{B} \mid \overline{A})=1-\mathrm{P}(B \mid \overline{A})=0{,}7$.
	Áp dụng công thức nhân xác suất, ta có xác suất trời nắng vào thứ Bảy và trời mưa vào Chủ nhật là
	$$
	\mathrm{P}(A B)=\mathrm{P}(A) \mathrm{P}(B \mid A)=0{,}7\cdot 0{,}2=0{,}14 .
	$$
Tương tự, ta có
	\allowdisplaybreaks
	\begin{eqnarray*}
		&&\mathrm{P}(A \overline{B})=\mathrm{P}(A) \mathrm{P}(\overline{B} \mid A)=0{,}7\cdot 0{,}8=0{,}56 ; \\
		&&\mathrm{P}(\overline{A} B)=\mathrm{P}(\overline{A}) \mathrm{P}(B \mid \overline{A})=0{,}3\cdot 0{,}3=0{,}09 ; \\
		&&\mathrm{P}(\overline{A} \overline{B})=\mathrm{P}(\overline{A}) \mathrm{P}(\overline{B} \mid \overline{A})=0{,}3\cdot0{,}7=0{,}21.
	\end{eqnarray*}
	Ta có thể biểu diễn các kết quả trên theo sơ đồ hình cây như sau:
	\begin{center}
		\begin{tikzpicture}
			\def\gocm{20}
			\def\gocn{10}
			\def\r{4}
			\tikzset{s/.style={outer sep=0.5 mm,draw=magenta,rectangle,minimum width=2.75cm,rounded corners=1mm}}
			\path(0,0)node(O){}++(\gocm:\r)node[s](A1){A}++(\gocn:\r)node[s](A2){$\overline{B}$}++(0:\r)node[s](A3){$A\overline{B}$}++(0:\r)node[s](A4){$0{,}56$};
			\path(A1)++({-\gocn}:\r)node[s](a2){B}++(0:\r)node[s](a3){$AB$}++(0:\r)node[s](a4){$0{,}14$};
			\path(O)++(-\gocm:\r)node[s](B1){$\overline{A}$}++(\gocn:\r)node[s](B2){$\overline{B}$}++(0:\r)node[s](B3){$\overline{A}\overline{B}$}++(0:\r)node[s](B4){$0{,}21$};
			\path(B1)++({-\gocn}:\r)node[s](b2){$B$}++(0:\r)node[s](b3){$\overline{A}B$}++(0:\r)node[s](b4){$0{,}09$};
			\foreach \x/\y in {
				O/A1,A1/A2,
				O/B1,B1/B2,
				A1/a2,
				B1/b2}
			\draw[-stealth](\x.east)--(\y.west);
			\path(O)--(A1.west)node[pos=0.5,above,sloped]{$0{,}7$}(O)--(B1.west)node[pos=0.5,below,sloped]{$0{,}3$}(B1.east)--(B2.west)node[pos=0.5,above,sloped]{$0{,}7$}(A1.east)--(A2.west)node[pos=0.5,above,sloped]{$0{,}8$}
			(A1.east)--(a2.west)node[pos=0.5,below,sloped]{$0{,}2$}
			(B1.east)--(b2.west)node[pos=0.5,below,sloped]{$0{,}3$};
			%%Node dòng trên
			\path(A2)++(0,1)node{\textbf{Chủ nhật}}++(180:4)node{\textbf{Thứ bảy}}(A3)++(0,1)node{\textbf{Kết quả}}(A4)++(0,1)node{\textbf{Xác suất}};
		\end{tikzpicture}
	\end{center}}
\end{bt}

\begin{bt}%[2D6H2-3]
% nguon: Archive/2024/TaiLieuGiangDay/6. K12 - 4in 1 - XsCoDieuKien/data/CHUONG 6 - XAC SUAT/Bai1-Dang34-SoDoCay-CongThucNhanXacSuat.tex #3256
%[BG12-4in1, Võ Thành Huấn]
	Hộp thứ nhất có $4$ viên bi xanh và $6$ viên bi đỏ. Hộp thứ hai có $5$ viên bi xanh và $4$ viên bi đỏ. Các viên bi có cùng kích thước và khối lượng. Lấy ra ngẫu nhiên $1$ viên bi từ hộp thứ nhất chuyển sang hộp thứ hai. Sau đó lại lấy ra ngẫu nhiên $1$ viên bi từ hộp thứ hai.
	Sử dụng sơ đồ hình cây, tính xác suất của các biến cố:
	$A$ : \lq\lq Viên bi lấy ra từ hộp thứ nhất có màu xanh và viên bi lấy ra từ hộp thứ hai có màu đỏ\rq\rq ;
	$B$ : \lq\lq Hai viên bi lấy ra có cùng màu\rq\rq.
	\loigiai{
		Gọi $X$ là biến cố: \lq\lq Viên bi lấy ra từ hộp thứ nhất có màu xanh\rq\rq.\\
		$Y$ là biến cố: \lq\lq Viên bi lấy ra từ hộp thứ hai có màu đỏ\rq\rq.\\
		Ta có
		$
		\mathrm{P}(Y|X)=0{,}4; \mathrm{P}(Y \mid \overline{X})=0{,}5 ; \mathrm{P}(X)=0{,}4
		$.\\
		Do đó $\mathrm{P}(\overline{X})=1-\mathrm{P}(X)=0{,}6 ; \mathrm{P}(\overline{Y} |X)=1-\mathrm{P}(Y|X)=0{,}6 ; \\\mathrm{P}(\overline{Y} \mid \overline{X})=1-\mathrm{P}(Y \mid \overline{X})=0{,}5$.\\
		Ta có sơ đồ hình cây như sau
		\begin{center}
			\begin{tikzpicture}
				\def\gocm{20}
				\def\gocn{10}
				\def\r{4}
				\tikzset{s/.style={outer sep=0.5 mm,draw=magenta,rectangle,minimum width=2.75cm,rounded corners=1mm}}
				\path(0,0)node(O){}++(\gocm:\r)node[s](A1){X}++(\gocn:\r)node[s](A2){$Y$};
				\path(A1)++({-\gocn}:\r)node[s](a2){$\overline{Y}$};
				\path(O)++(-\gocm:\r)node[s](B1){$\overline{X}$}++(\gocn:\r)node[s](B2){$Y$};
				\path(B1)++({-\gocn}:\r)node[s](b2){$\overline{Y}$};
				\foreach \x/\y in {
					O/A1,A1/A2,
					O/B1,B1/B2,
					A1/a2,
					B1/b2}
				\draw[-stealth](\x.east)--(\y.west);
				\path(O)--(A1.west)node[pos=0.5,above,sloped]{$0{,}4$}(O)--(B1.west)node[pos=0.5,below,sloped]{$0{,}6$}(B1.east)--(B2.west)node[pos=0.5,above,sloped]{$0{,}5$}(A1.east)--(A2.west)node[pos=0.5,above,sloped]{$0{,}4$}
				(A1.east)--(a2.west)node[pos=0.5,below,sloped]{$0{,}6$}
				(B1.east)--(b2.west)node[pos=0.5,below,sloped]{$0{,}5$};
				
			\end{tikzpicture}
			
		\end{center}
		Khi đó $\mathrm{P}(A)=\mathrm{P}(XY)=0{,}4\cdot 0{,}4=0{,}16$;\\ $\mathrm{P}(B)=\mathrm{P}(X\overline{Y})+\mathrm{P}(\overline{X}Y)=0{,}4\cdot 0{,}6+0{,}6\cdot 0{,}5=0{,}54$.
	}
\end{bt}

\begin{bt}%[2D6V2-3]
% nguon: Archive/2024/TaiLieuGiangDay/6. K12 - 4in 1 - XsCoDieuKien/data/CHUONG 6 - XAC SUAT/Bai1-Dang34-SoDoCay-CongThucNhanXacSuat.tex #3257
%[BG12-4in1, Võ Thành Huấn]
	Một trường đại học tiến hành khảo sát tình trạng việc làm sau khi tốt nghiệp của sinh viên. Kết quả khảo sát cho thấy tỉ lệ người tìm được việc làm đúng chuyên ngành là $85 \%$ đối với sinh viên tốt nghiệp loại giòi và $70 \%$ đối với sinh viên tốt nghiệp loại khác.
	Tỉ lệ sinh viên tốt nghiệp loại giòi là $30 \%$. Gặp ngẫu nhiên một sinh viên đã tốt nghiệp của trường.
	Sử dụng sơ đồ hình cây, tính xác suất của các biến cố:
	C: \lq\lq Sinh viên tốt nghiệp loại giỏi và tìm được việc làm đúng chuyên ngành\rq\rq;
	$D$ : \lq\lq Sinh viên không tốt nghiệp loại giỏi và tìm được việc làm đúng chuyên ngành\rq\rq.
	\loigiai{
		Gọi $X$ là biến cố: \lq\lq Sinh viên tốt nghiệp loại Giỏi\rq\rq.\\
		$Y$ là biến cố: \lq\lq Sinh viên tìm được việc làm đúng chuyên ngành\rq\rq.\\
		Ta có
		$
		\mathrm{P}(Y|X)=0{,}85; \mathrm{P}(Y \mid \overline{X})=0{,}7 ; \mathrm{P}(X)=0{,}3
		$.\\
		Do đó $\mathrm{P}(\overline{X})=1-\mathrm{P}(X)=0{,}7 ; \mathrm{P}(\overline{Y} |X)=1-\mathrm{P}(Y|X)=0{,}15 ; \\\mathrm{P}(\overline{Y} \mid \overline{X})=1-\mathrm{P}(Y \mid \overline{X})=0{,}3$.\\
		Ta có sơ đồ hình cây như sau
		\begin{center}
			\begin{tikzpicture}
				\def\gocm{20}
				\def\gocn{10}
				\def\r{4}
				\tikzset{s/.style={outer sep=0.5 mm,draw=magenta,rectangle,minimum width=2.75cm,rounded corners=1mm}}
				\path(0,0)node(O){}++(\gocm:\r)node[s](A1){X}++(\gocn:\r)node[s](A2){$Y$};
				\path(A1)++({-\gocn}:\r)node[s](a2){$\overline{Y}$};
				\path(O)++(-\gocm:\r)node[s](B1){$\overline{X}$}++(\gocn:\r)node[s](B2){$Y$};
				\path(B1)++({-\gocn}:\r)node[s](b2){$\overline{Y}$};
				\foreach \x/\y in {
					O/A1,A1/A2,
					O/B1,B1/B2,
					A1/a2,
					B1/b2}
				\draw[-stealth](\x.east)--(\y.west);
				\path(O)--(A1.west)node[pos=0.5,above,sloped]{$0{,}3$}(O)--(B1.west)node[pos=0.5,below,sloped]{$0{,}7$}(B1.east)--(B2.west)node[pos=0.5,above,sloped]{$0{,}7$}(A1.east)--(A2.west)node[pos=0.5,above,sloped]{$0{,}85$}
				(A1.east)--(a2.west)node[pos=0.5,below,sloped]{$0{,}15$}
				(B1.east)--(b2.west)node[pos=0.5,below,sloped]{$0{,}3$};
				
			\end{tikzpicture}
		\end{center}
		Khi đó $\mathrm{P}(C)=\mathrm{P}(XY)=0{,}3\cdot 0{,}85=0{,}255$; $\mathrm{P}(B)=\mathrm{P}(\overline{X}Y)=0{,}7\cdot 0{,}7=0{,}49$.
	}
\end{bt}

\begin{bt}%[2D6H2-4]
% nguon: Archive/2024/TaiLieuGiangDay/6. K12 - 4in 1 - XsCoDieuKien/data/CHUONG 6 - XAC SUAT/Bai1-Dang34-SoDoCay-CongThucNhanXacSuat.tex #3258
[KNTT]%[BG12-4in1, Võ Thành Huấn]
	Trong một hộp kín có 7 chiếc bút bi xanh và 5 chiếc bút bi đen, các chiếc bút có cùng kích thước và khối lượng. Bạn Sơn lấy ngẫu nhiên một chiếc bút bi từ trong hộp, không trả lại. Sau đó bạn Tùng lấy ngẫu nhiên một trong 11 chiếc bút còn lại. Tính xác suất để Sơn lấy được bút bi đen và Tùng lấy được bút bi xanh.
	\loigiai{
		Gọi $A$ là biến cố: ``Bạn Sơn lấy được bút bi đen'';\\
		$B$ là biến cố: ``Bạn Tùng lấy được bút bi xanh''.\\
		Ta cần tính $\mathrm{P}(AB)$.\\
		Vì $n(A)=5$ nên $\mathrm{P}(A)=\dfrac{5}{12}$.\\
		Nếu $A$ xảy ra tức là bạn Sơn lấy được bút bi đen thì trong hộp có 11 bút bi với 7 bút bi xanh.\\
		Vậy $\mathrm{P}(B\mid A)=\dfrac{7}{11}$.\\
		Theo công thức nhân xác suất: $\mathrm{P}(AB)=\mathrm{P}(A)\cdot \mathrm{P}(B\mid A)=\dfrac{5}{12}\cdot\dfrac{7}{11}=\dfrac{35}{132}$.\\}
\end{bt}

\begin{bt}%[2D6H2-4]
% nguon: Archive/2024/TaiLieuGiangDay/6. K12 - 4in 1 - XsCoDieuKien/data/CHUONG 6 - XAC SUAT/Bai1-Dang34-SoDoCay-CongThucNhanXacSuat.tex #3259
[KNTT]%[BG12-4in1, Võ Thành Huấn]
	Ô cửa bí mật (Let's Make a Deal) là một trò chơi trên truyền hình nổi tiếng ở Mỹ, đã được mua bản quyền và phát sóng ở nhiều nước trên thế giới. Nội dung trò chơi như sau:
	\begin{enumerate}[$\bullet$]
		\item Người chơi được mời lên sân khấu và đứng trước ba cánh cửa đóng kín. Sau một cánh cửa có chiếc ô tô, sau mỗi cánh cửa còn lại là một con lừa. Người chơi được yêu cầu chọn ngẫu nhiên một cánh cửa, nhưng không được mở ra.
		\item Tiếp đó người quản trò tuyên bố sẽ mở ngẫu nhiên một trong hai cánh cửa người chơi không chọn mà sau cửa đó là con lừa. Người quản trò hỏi người chơi muốn giữ nguyên sự lựa chọn ban đầu của mình hay muốn chuyển sang cửa chưa mở còn lại.
	\end{enumerate}
	Giả sử người chơi chọn cửa số 1 và người quản trò mở cửa số 3. Kí hiệu $E_1$; $E_2$; $E_3$ tương ứng là các biến cố: ``Sau ở cửa số 1 có ô tô''; ``Sau ở cửa số 2 có ô tô''; ``Sau ở cửa số 3 có ô tô'' và $H$ là biến cố: ``Người quản trò mở ở cửa số 3 thấy con lừa''.
	Sau khi người quản trò mở cánh cửa số 3 thấy con lừa, tức là khi $H$ xảy ra. Để quyết định thay đổi lựa chọn hay không, người chơi cần so sánh hai xác suất có điều kiện: $P\left(E_1\mid H\right)$ và $P\left(E_2\mid H\right)$
	\begin{enumerate}[a)]
		\item Chứng minh rằng:
		\begin{enumerate}[$\bullet$]
			\item $\mathrm{P}(E_1)=\mathrm{P}(E_2)=\mathrm{P}(E_3)=\dfrac{1}{3}$;
			\item $P\left(H \mid E_1\right)=\dfrac{1}{2} ~\text{và}~P\left(H \mid E_2\right)=1$.
		\end{enumerate}
		\item Sử dụng công thức tính xác suất có điều kiện và công thức nhân xác suất, chứng minh rằng:
		\begin{enumerate}[$\bullet$]
			\item $ \mathrm{P}\left(E_1 \mid H\right)=\dfrac{\mathrm{P}(E_1) \cdot P\left(H \mid E_1\right)}{\mathrm{P}(H)}$;
			%		\item $P\left(E_2 \mid H\right)=\dfrac{\mathrm{P}(E_2) \cdot P\left(H \mid E_2\right)}{\mathrm{P}(H)$.
			\end{enumerate}
			\item Từ các kết quả trên hãy suy ra: 
			$$P\left(E_2 \mid H\right)=2 P\left(E_1 \mid H\right)$$
			Từ đó hãy đưa ra lời khuyên cho người chơi: Nên giữ nguyên sự lựa chọn ban đầu hay chuyển sang cửa chưa mở còn lại?\\
		\end{enumerate}
		\loigiai{
			\begin{enumerate}[a)]
				\item Không gian mẫu $\Omega$ là tập hợp gồm 3 phần thưởng (1 ô tô + 2 con lừa) $\Rightarrow n(\Omega)=3$.
				Ta có $n(E_1)=n(E_2)=n(E_3)=1$.\\
				Suy ra $\mathrm{P}(E_1)=\mathrm{P}(E_2)=\mathrm{P}(E_3)=\dfrac{1}{3}$.\\
				Nếu $E_1$ xảy ra, tức là sau cửa số 1 có ô tô. Khi đó, sau cửa số 2 và 3 là con lừa. Người quản trò chọn ngẫu nhiên một trong hai cửa số 2 và 3 để mở ra. Do đó, việc chọn cửa số 2 hay cửa số 3 có khả năng như nhau. \\
				Vậy $P\left(H\mid E_1\right)=\dfrac{1}{2}$.\\
				Nếu $E_2$ xảy ra, tức là cửa số 2 có ô tô. Khi đó, người quản trò chắc chắn phải mở cửa số 3.\\
				Do đó $P\left(H\mid E_2\right)=1$.
				\item Ta có $$\begin{aligned} &P\left(E_1\mid H\right)=\dfrac{\mathrm{P}(E_1H)}{\mathrm{P}(H)}\\\Leftrightarrow &\mathrm{P}(E_1H)=\mathrm{P}(E_1)\cdot P\left(H\mid E_1\right)\\\Leftrightarrow &P\left(E_1\mid H\right)=\dfrac{\mathrm{P}(E_1)\cdot P\left(H\mid E_1\right)}{\mathrm{P}(H)}.\end{aligned}$$
				\item Ta có $$P\left(E_2 \mid H\right)=\dfrac{\mathrm{P}(E_2) \cdot P\left(H\mid E_2\right)}{\mathrm{P}(H)}.$$
				Suy ra 
				$$\begin{aligned}
					\dfrac{P\left(E_2 \mid H\right)}{P\left(E_1 \mid H\right)}&=\dfrac{\mathrm{P}(E_2) \cdot P\left(H \mid E_2\right)}{\mathrm{P}(E_1) \cdot P\left(H \mid E_1\right)}\\
					&=\dfrac{\dfrac{1}{3} \cdot 1}{\dfrac{1}{3} \cdot\dfrac{1}{2}}=2.
				\end{aligned}$$
				Suy ra $P\left(E_2\mid H\right)=2\cdot P\left(E_1\mid H\right)$.\\
				\textbf{\textit{Nhận xét: }}Từ kết quả trên ta thấy người chơi nên chuyển sang cửa chưa mở còn lại để tăng gấp đôi khả năng trúng thưởng chiếc ô tô.
			\end{enumerate}
		}
\end{bt}

\begin{bt}%[2D6H1-2]
% nguon: Archive/2024/TaiLieuGiangDay/6. K12 - 4in 1 - XsCoDieuKien/data/CHUONG 6 - XAC SUAT/Bai1-Dang34-SoDoCay-CongThucNhanXacSuat.tex #3261
%[BG12-4in1, Võ Thành Huấn]
	Cho hai biến cố $A$ và $B$ có $\mathrm{P}(A)=0{,}4; \mathrm{P}(B)=0{,}8$ và $\mathrm{P}(A|B)=0{,}5$. Tính $\mathrm{P}(A\overline{B})$ và $\mathrm{P}(A|B)$.
	\loigiai{
		Ta có $\mathrm{P}(A\overline{B})=\mathrm{P}(A|\overline{B})\cdot \mathrm{P}(\overline{B})=0{,}5\cdot 0{,}2=0{,}1$.\\
		Vì $AB$ và $A\overline{B}$ là hai biến cố xung khắc và $AB\cup A\overline{B}=A$ nên theo tính chất của xác suất, ta có\\ $\mathrm{P}(AB)=\mathrm{P}(A)-\mathrm{P}(A\overline{B})=0{,}4-0{,}1=0{,}3$.\\
		Khi đó: $\mathrm{P}(A|B)=\dfrac{\mathrm{P}(AB)}{\mathrm{P}(B)}=\dfrac{0{,}3}{0{,}8}=0{,}375$.
	}
\end{bt}

\begin{bt}%[2D6H2-4]
% nguon: Archive/2024/TaiLieuGiangDay/6. K12 - 4in 1 - XsCoDieuKien/data/CHUONG 6 - XAC SUAT/Bai1-Dang34-SoDoCay-CongThucNhanXacSuat.tex #3263
[CD]%[BG12-4in1, Võ Thành Huấn]
	Một lô sản phẩm có $20$ sản phẩm, trong đó có $5$ sản phẩm chất lượng thấp. Lấy liên tiếp $2$ sản phẩm trong lô sản phẩm trên, trong đó sản phẩm lấy ra ở lần thứ nhất không được bỏ lại vào lô sản phẩm. Tính xác suất để cả hai sản phẩm được lấy ra đều có chất lượng thấp.
	\loigiai{
		Gọi $A$ là biến cố \lq\lq sản phẩm thứ nhất có chất lượng thấp\rq\rq, và $B$ là biến cố \lq\lq sản phẩm thứ hai có chất lượng thấp\rq\rq.\\
		Xác suất của $A$ là xác suất để lấy ra một sản phẩm chất lượng thấp trong lần đầu tiên:
		$$\mathrm{P}(A)=\dfrac{n(A)}{n(\Omega)}=\dfrac{5}{20}=\dfrac{1}{4}$$
		Sau khi lấy một sản phẩm chất lượng thấp, số sản phẩm chất lượng thấp giảm còn $4$ trong tổng số $19$ sản phẩm.\\
		Xác suất của $B$ khi đã xảy ra $A$ là xác suất để lấy ra một sản phẩm chất lượng thấp trong lần thứ hai:
		$$\mathrm{P}(B \mid A)=\dfrac{4}{19}$$
		Áp dụng quy tắc nhân xác suất:
		$$\mathrm{P}(A \cap B)=\mathrm{P}(A)\cdot \mathrm{P}(B \mid A)=\dfrac{1}{4}\cdot\dfrac{4}{19}=\dfrac{1}{19}$$
		Vậy, xác suất để cả hai sản phẩm được lấy ra đều có chất lượng thấp là $\dfrac{1}{19}$.
	}
\end{bt}

\begin{bt}%[2D6H2-2]
% nguon: Archive/2024/TaiLieuGiangDay/6. K12 - 4in 1 - XsCoDieuKien/data/CHUONG 6 - XAC SUAT/Bai2-Dang1-CongThucXSTP.tex #3271
[CD, Mức độ 1]%[Trần Ngọc Thành]
	Cho hai biến cố $A$, $B$ với $\mathrm{P}(B)=0{,}6$; $\mathrm{P}(A|B) =0{,}7$ và $\mathrm{P}\left(A|\overline{B}\right)=0{,}4$. Tính $\mathrm{P}(A)$
	\loigiai{
		Ta có $\mathrm{P}\left(\overline{B}\right)= 1- \mathrm{P}(B) = 1-0{,}6 = 0{,}4$.\\
		Áp dụng công thức xác suất toàn phần, ta có
		\[\mathrm{P}(A) = \mathrm{P}(A|B)\cdot \mathrm{P}(B) + \mathrm{P}\left(A|\overline{B}\right)\cdot \mathrm{P}\left(\overline{B}\right)= 0{,}7\cdot 0{,}6 + 0{,}4\cdot 0{,}4=0{,}58.\]
	}
\end{bt}

\begin{bt}%[2D6H2-2]
% nguon: Archive/2024/TaiLieuGiangDay/6. K12 - 4in 1 - XsCoDieuKien/data/CHUONG 6 - XAC SUAT/Bai2-Dang1-CongThucXSTP.tex #3273
[KNTT, Mức độ 2]%[Trần Ngọc Thành]
	Trong một kì thi tốt nghiệp trung học phổ thông, một tỉnh X có $80 \%$ học sinh lựa chọn tổ hợp A00 (gồm các môn Toán, Vật lí, Hoá học). Biết rằng, nếu một học sinh chọn tổ hợp A00 thì xác suất để học sinh đó đỗ đại học là 0{,}6; còn nếu một học sinh không chọn tổ hợp A00 thì xác suất để học sinh đó đỗ đại học là 0{,}7. Chọn ngẫu nhiên một học sinh của tỉnh X đã tốt nghiệp trung học phổ thông trong kì thi trên. Tính xác suất để học sinh đó đỗ đại học.
	\loigiai{Gọi $A$ là biến cố: \lq \lq Học sinh đó chọn tổ hợp A00\rq \rq~; $B$ là biến cố: \lq \lq Học sinh đó đỗ đại học\rq \rq.\\
		Ta cần tính $\mathrm{P}(B)$. Theo công thức xác suất toàn phần, ta cần biết: $\mathrm{P}(A), \mathrm{P}(\overline{A}), \mathrm{P}(B \mid A)$ và $\mathrm{P}(B \mid \overline{A})$.\\
		Ta có: $\mathrm{P}(A)=0{,}8; \mathrm{P}(\overline{A})=1-\mathrm{P}(A)=1-0{,}8=0{,}2$.\\
		$\mathrm{P}(B \mid A)$ là xác suất để một học sinh đỗ đại học với điều kiện học sinh đó chọn tổ hợp $A 00$ \\$\Rightarrow \mathrm{P}(B \mid A)=0{,}6$.\\
		$\mathrm{P}(B \mid \overline{A})$ là xác suất để một học sinh đỗ đại học với điều kiện học sinh đó không chọn tổ hợp $\mathrm{A} 00$\\$ \Rightarrow \mathrm{P}(B \mid \overline{A})=0{,}7$.\\
		Thay vào công thức xác suất toàn phần ta được:$$\mathrm{P}( B)={\mathrm{P}(A) \cdot \mathrm{P}(B \mid A)+\mathrm{P}(\overline{A}) \cdot \mathrm{P}(B \mid \overline{A})}={0{,}8 \cdot 0{,}6+0{,}2 \cdot 0{,}7} = 0{,}62.$$}
\end{bt}

\begin{bt}%[2D6H2-2]
% nguon: Archive/2024/TaiLieuGiangDay/6. K12 - 4in 1 - XsCoDieuKien/data/CHUONG 6 - XAC SUAT/Bai2-Dang1-CongThucXSTP.tex #3274
[CD, Mức độ 3]%[Trần Ngọc Thành]
	Một hộp có $60$ viên bi màu xanh và $40$ viên bi màu đỏ; các viên bi có kích thước và khối lượng như nhau. Sau khi thống kê, người ta thấy: có $50\%$ số viên bi màu xanh có dán nhãn và $75\%$ số viên bi màu đỏ có dán nhãn; những viên bi còn lại không có dán nhãn.
	\begin{enumerate}
		\item Chọn số thích hợp cho $\boxed{?}$ trong bảng (đơn vị: viên bi).
		\begin{center}
			\begin{tabular}{|c|C{3cm}|C{3cm}|}
				\hline 
				\diaghead{Dán nhãn Màu bi}{\normalsize Màu bi}{\normalsize Dán nhãn} & Có dán nhãn & Không dán nhãn \\
				\hline 
				Đỏ & $\boxed{?}$ & $\boxed{?}$ \\
				\hline
				Xanh & $\boxed{?}$ & $\boxed{?}$\\
				\hline
			\end{tabular}
		\end{center}
		\item Lấy ra ngẫu nhiên một viên bi trong hộp. Sử dụng công thức xác suất toàn phần, tính xác suất để viên bi được lấy ra có dán nhãn.
	\end{enumerate}
	\loigiai{
		\begin{enumerate}
			\item Số viên bi màu đỏ có dán nhãn là $75\%\cdot 40 = 30$ (viên bi).\\
			Số viên bi màu xanh có dãn nhẫn là $50\%\cdot 60 = 30$ (viên bi).\\
			\begin{center}
				\begin{tabular}{|c|C{2.75cm}|C{2.75cm}|}
					\hline 
					\diaghead{Dán nhãn Màu bi}{\normalsize Dán nhãn}{\normalsize Màu bi} & Có dán nhãn & Không dán nhãn \\
					\hline 
					Đỏ & $30$ & $10$ \\
					\hline
					Xanh & $30$ & $30$\\
					\hline
				\end{tabular}
			\end{center}
			Sau khi hoàn thiện bảng $3$ ta được bảng $4$ (đơn vị: viên bi).
			\item Xét hai biến cố sau
			\begin{itemize}
				\item[] $A$: \lq\lq Viên bi được chọn ra có dãn nhãn\rq\rq.
				\item[] $B$: \lq\lq Viên bi được chọn ra có màu đỏ\rq\rq. 
			\end{itemize} 
			Khi đó, ta có
			\begin{listEX}[3]
				\item[] $\mathrm{P}(B)=\dfrac{40}{100} = \dfrac{2}{5}$;
				\item*[] $\mathrm{P}\left(\overline{B}\right)= 1 - \mathrm{P}(B)=1-\dfrac{2}{5}=\dfrac{3}{5}$;
				\item[] $\mathrm{P}(A|B) = \dfrac{30}{40}=\dfrac{3}{4}$;
				\item*[] $\mathrm{P}\left(A|\overline{B}\right)=\dfrac{30}{60}=\dfrac{1}{2}$.
			\end{listEX}
			Áp dụng công thức tính xác suất toàn phần, ta có 
			\[\mathrm{P}(A) = \mathrm{P}(B)\cdot\mathrm{P}\left(A|B\right) + \mathrm{P}\left(\overline{B}\right)\cdot\mathrm{P}\left(A|\overline{B}\right) = \dfrac{2}{3}\cdot \dfrac{3}{4} + \dfrac{3}{5}\cdot \dfrac{1}{2}=\dfrac{4}{5}.\]
			Vậy xác suất để viên bi được lấy ra có dãn nhãn bằng $\dfrac{4}{5}$.
		\end{enumerate}
	}
\end{bt}

\begin{bt}%[2D6H2-2]
% nguon: Archive/2024/TaiLieuGiangDay/6. K12 - 4in 1 - XsCoDieuKien/data/CHUONG 6 - XAC SUAT/Bai2-Dang1-CongThucXSTP.tex #3275
[CD, Mức độ 3]%[Trần Ngọc Thành]
	Trong trò chơi hái hoa có thưởng của lớp 12A, cô giáo treo $10$ bông hoa trên cành cây, trong đó có $5$ bông hoa chưa phiếu có thưởng. Bạn Bình hái bông hoa đầu tiên, sau đó bạn An hái bông hoa thứ hai.
	\begin{enumerate}
		\item Vẽ sơ đồ cây biểu thị tình huống trên.
		\item Từ đó, tính xác suất bạn An hái được bông hoa chứa phiếu có thưởng.
	\end{enumerate}
	\loigiai{
		Xét hai biến cố
		\begin{itemize}
			\item[] $A$: \lq\lq Bông hoa bạn An hái được chứa phiếu có thưởng\rq\rq.
			\item[] $B$: \lq\lq Bông hoa bạn Bình hái được chứa phiếu có thưởng\rq\rq. 
		\end{itemize}
		Khi đó, ta có
		\begin{listEX}[3]
			\item[] $\mathrm{P}(B)=\dfrac{5}{10}=\dfrac{1}{2}$
			\item*[] $\mathrm{P}\left(\overline{B}\right) = 1 - \mathrm{P}(B)=1-\dfrac{1}{2}=\dfrac{1}{2}$
			\item[] $\mathrm{P}(A|B)=\dfrac{4}{9}$
			\item*[] $\mathrm{P}\left(A|\overline{B}\right)=\dfrac{5}{9}$.
		\end{listEX}
		\begin{enumerate}
			\item Sơ đồ hình cây biểu thị tình huống đã cho là
			\begin{center}
				\begin{tikzpicture}[->,>=stealth,line join=round,line cap=round,font=\footnotesize,scale=1]
					\def\xmot{3}
					\def\xhai{8}
					\node (O) at (0,0){};
					\node (B) at (\xmot,1){$B$};
					\node (B1) at (\xmot,-1){$\overline{B}$};
					\node (BA) at (\xhai,2){$A$};
					\node (BA1) at (\xhai,0.3){$\overline{A}$};
					\node (B1A) at (\xhai,-0.3){$A$};
					\node (B1A1) at (\xhai,-1.75){$\overline{A}$};
					
					\foreach \x/\y/\p/\l in 
					{
						O/B/above/$\mathrm{P}(B)=\dfrac{1}{2}$,
						B/BA/above/$\mathrm{P}(A|B)=\dfrac{4}{9}$,
						B/BA1//,
						O/B1/below/$\mathrm{P}\left(\overline{B}\right)=\dfrac{1}{2}$,
						B1/B1A/above/$\mathrm{P}\left(A|\overline{B}\right)=\dfrac{5}{9}$,
						B1/B1A1//
					}
					{
						\draw[->] (\x)--(\y)node[midway,\p,scale=0.8,sloped]{\l};
					}
					\node (mot) at (\xmot,4) {Bông hoa bạn Bình hái ra};
					\node (hai) at (\xhai,4) {Bông hoa bạn An hái ra};
					\draw[->,thick,orange] (mot)--(B);
					\draw[->,thick,orange] (hai)--(BA);
				\end{tikzpicture}
			\end{center}
			\item Áp dụng công thức tính xác suất toàn phần, ta có:
			\[\mathrm{P}(A) = \mathrm{P}(B)\cdot\mathrm{P}\left(A|B\right) + \mathrm{P}\left(\overline{B}\right)\cdot\mathrm{P}\left(A|\overline{B}\right) = \dfrac{1}{2}\cdot \dfrac{4}{9} + \dfrac{1}{2}\cdot\dfrac{5}{9}=\dfrac{1}{2}.\]
			Vậy xác suất bạn An hái được bông hoa chứa phiếu có thưởng bằng $\dfrac{1}{2}$.
		\end{enumerate}
	}
\end{bt}

\begin{bt}%[2D6H2-2]
% nguon: Archive/2024/TaiLieuGiangDay/6. K12 - 4in 1 - XsCoDieuKien/data/CHUONG 6 - XAC SUAT/Bai2-Dang1-CongThucXSTP.tex #3278
[CD, Mức độ 2]%[Trần Ngọc Thành]
	Một loại linh kiện do hai nhà máy số I, số II cùng sản xuất. Tỉ lệ phế phẩm của các nhà máy I, II lần lượt là $4\%$; $3\%$. Trong một lô linh kiện để lẫn lộn $80$ sản phẩm của nhà máy số I và $120$ sản phẩm của nhà máy số II. Một khách hàng lấy ngẫu nhiên một linh liện từ lô hàng đó.
	Tính xác suất để linh kiện được lấy ra là linh kiện tốt.
	\loigiai{
		
		Xét các biến cố
		\begin{itemize}
			\item[] $A$: \lq\lq Linh kiện lấy ra là linh kiện tốt\rq\rq.
			\item[] $B$: \lq\lq Linh kiện lấy ra là linh kiện từ nhà máy số I\rq\rq.
			\item[] $\overline{B}$: \lq\lq Linh kiện lấy ra là linh kiện từ nhà máy số II\rq\rq.
		\end{itemize}
		Theo đề bài, ta có
		\begin{listEX}[2]
			\item[] $\mathrm{P}(A|B) =1 - 0{,}04 = 0{,}96$.
			\item[] $\mathrm{P}\left(A|\overline{B}\right) = 1-0{,}03 = 0{,}97$.
			\item[] $\mathrm{P}(B)=\dfrac{80}{200}=0{,}4$.
			\item[] $\mathrm{P}\left(\overline{B}\right) = \dfrac{120}{200}=0{,}6$.
		\end{listEX}
		Khi đó áp dụng công thức xác suất toàn phần, ta có
		\[\mathrm{P}(A) = \mathrm{P}(A|B)\cdot \mathrm{P}(B) + \mathrm{P}\left(A|\overline{B}\right)\cdot\mathrm{P}\left(\overline{B}\right)=0{,}96\cdot 0{,}4 + 0{,}97\cdot 0{,}6=0{,}966.\]
		
		
	}
\end{bt}

\begin{bt}%[2D6H2-2]
% nguon: Archive/2024/TaiLieuGiangDay/6. K12 - 4in 1 - XsCoDieuKien/data/CHUONG 6 - XAC SUAT/Bai2-Dang1-CongThucXSTP.tex #3279
[CKP, Mức độ 2]%[Trần Ngọc Thành]
	Người ta khảo sát khả năng chơi nhạc cụ của một nhóm học sinh tại trường X. Nhóm này có $60\%$ học sinh là nam. Kết quả khảo sát cho thấy có $20\%$ học sinh nam và $15\%$ học sinh nữ biết chơi ít nhất một nhạc cụ. Chọn ngẫu nhiên một học sinh trong nhóm này. Tính xác suất để chọn được học sinh biết chơi ít nhất một nhạc cụ.
	\loigiai{
		Xét phép thử chọn ngẫu nhiên một học sinh trong nhóm.\\
		Gọi $A$ là biến cố \lq \lq Chọn được một học sinh biết chơi ít nhất một nhạc cụ\rq \rq~ và $B$, $\overline{B}$ lần lượt là các biến cố \lq \lq Chọn được một học sinh nam\rq \rq~ và \lq \lq Chọn được một học sinh nữ\rq \rq~.\\
		Theo đề bài: 
		$$\mathrm{P}(B) = 60\% = 0{,}6;\quad \mathrm{P}(\overline{B}) = 1 - 0{,}6 = 0{,}4;$$
		$$\mathrm{P}(A|B) = 20\% = 0{,}2;\quad \mathrm{P}(A|\overline{B}) = 15\% = 0{,}15.$$
		Áp dụng công thức xác suất toàn phần, ta có:
		$$\mathrm{P}(A) = \mathrm{P}(B)\cdot \mathrm{P}(A|B) + \mathrm{P}(\overline{B})\cdot \mathrm{P}(A|\overline{B}) = 0{,}6\cdot 0{,}2 + 0{,}4\cdot 0{,}15 = 0{,}18.$$
		Vậy xác suất để chọn được một học sinh biết chơi nhạc cụ là $0{,}18$.
	}
\end{bt}

\begin{bt}%[2D6H2-2]
% nguon: Archive/2024/TaiLieuGiangDay/6. K12 - 4in 1 - XsCoDieuKien/data/CHUONG 6 - XAC SUAT/Bai2-Dang1-CongThucXSTP.tex #3281
[KNTT, Mức độ 2]%[Trần Ngọc Thành]
	Có hai đội thi đấu môn Bắn súng. Đội I có 5 vận động viên, đội II có 7 vận động viên. Xác suất đạt huy chương vàng của mỗi vận động viên đội I và đội II tương ứng là 0{,}65 và 0{,}55. Chọn ngẫu nhiên một vận động viên.
	Tính xác suất để vận động viên này đạt huy chương vàng;
	\loigiai{
		Xét biến cố $A$: \lq \lq Vận động viên này thuộc đội I\rq \rq~. Xét biến cố $B$: \lq \lq Vận động viên này đạt huy chương vàng\rq \rq~.
		Ta có $\mathrm{P}(B)=\mathrm{P}(A)\cdot \mathrm{P}(B \mid A)+\mathrm{P}(\overline{A}) \cdot \mathrm{P}(B \mid \overline{A})$.
		\begin{itemize}
			\item Tính $\mathrm{P}(A)$: Đây là xác suất để vận động viên đó thuộc đội I. Vậy $\mathrm{P}(A)=\dfrac{5}{12}$.
			\item Tính $\mathrm{P}(\overline{A})$: $\mathrm{P}(\overline{A})=1-\mathrm{P}(A)=\dfrac{7}{12}$.
			\item Tính $\mathrm{P}(B\mid A)$: Đây là xác suất để vận động viên thuộc đội I đạt huy chương vàng.\\ Vậy $\mathrm{P}(B\mid A)=0{,}65$.
			\item Tính $\mathrm{P}(B\mid \overline{A})$: Đây là xác suất để vận động viên thuộc đội I đạt huy chương vàng. \\ Vậy $\mathrm{P}(B\mid \overline{A})=0{,}55$.
		\end{itemize}
		Vậy $\mathrm{P}(B)=\mathrm{P}(A)\cdot \mathrm{P}(B \mid A)+\mathrm{P}(\overline{A}) \cdot \mathrm{P}(B \mid \overline{A})=\dfrac{5}{12}\cdot 0{,}65+\dfrac{7}{12}\cdot 0{,}55=\dfrac{71}{120}\approx 0{,}59$. \\
		Vậy xác suất để vận động viên này đạt huy chương vàng là khoảng $0{,}59$.
	}
\end{bt}

\begin{bt}%[2D6H2-2]
% nguon: Archive/2024/TaiLieuGiangDay/6. K12 - 4in 1 - XsCoDieuKien/data/CHUONG 6 - XAC SUAT/Bai2-Dang1-CongThucXSTP.tex #3282
[KNTT, Mức độ 3]%[Trần Ngọc Thành]
	Chuồng I có 5 con gà mái, 2 con gà trống. Chuồng II có 3 con gà mái, 5 con gà trống. Bác Mai bắt một con gà trong số đó theo cách sau: Bác tung một con xúc xắc cân đối, đồng chất. Nếu số chấm chia hết cho 3 thì bác chọn chuồng I, nếu số chấm không chia hết cho 3 thì bác chọn chuồng II. Sau đó, từ chuồng đã chọn bác bắt ngẫu nhiên một con gà. Tính xác suất để bác Mai bắt được con gà mái.
	\loigiai{
		Gọi $A$ là biến cố: \lq\lq Bác Mai bắt được con gà mái\rq\rq.\\
		Gọi $B$ là biến cố: \lq\lq tung con xúc xắc được số chấm chia hết cho $3$\rq\rq\, suy ra $\mathrm{P}(B)=\dfrac{2}{6}=\dfrac{1}{3}$.\\
		Gọi $\overline{B}$ là biến cố: \lq\lq tung con xúc xắc được số chấm không chia hết cho $3$\rq\rq\, suy ra $\mathrm{P}(\overline{B})=\dfrac{4}{6}=\dfrac{2}{3}$.\\
		Ta có  xác suất bắt được gà mái từ chuồng I là $\mathrm{P}\left(A\mid B\right)=\dfrac{5}{7}$.\\
		Ta có  xác suất bắt được gà mái từ chuồng II là $\mathrm{P}\left(A\mid \overline{B}\right)=\dfrac{3}{8}$.\\
		Áp dụng công thức xác suất toàn phần ta có
		$$\mathrm{P}(A)=\mathrm{P}(B)\cdot \mathrm{P}(A\mid B)+\mathrm{P}(\overline{B})\cdot \mathrm{P}(A\mid \overline{B})=\dfrac{1}{3}\cdot \dfrac{5}{7}+\dfrac{2}{3}\cdot \dfrac{3}{8}=\dfrac{41}{84}.$$
		Vậy xác suất để bác Mai bắt được con gà mái là $\dfrac{41}{84}$.
	}
\end{bt}

\begin{bt}%[2D6H2-2]
% nguon: Archive/2024/TaiLieuGiangDay/6. K12 - 4in 1 - XsCoDieuKien/data/CHUONG 6 - XAC SUAT/Bai2-Dang1-CongThucXSTP.tex #3283
[KNTT, Mức độ 3]%[Trần Ngọc Thành]
	Tại nhà máy X sản xuất linh kiện điện tử tỉ lệ sản phẩm đạt tiêu chuẩn là $80 \%$. Trước khi xuất xưởng ra thị trường, các linh kiện điện tử đều phải qua khâu kiểm tra chất lượng để đóng dấu OTK. Vì sự kiểm tra không tuyệt đối hoàn hảo nên nếu một linh kiện điện tử đạt tiêu chuẩn thì nó có xác suất 0{,}99 được đóng dấu OTK; nếu một linh kiện điện tử không đạt tiêu chuẩn thì nó có xác suất 0{,}95 không được đóng dấu OTK. Chọn ngẫu nhiên một linh kiện điện tử của nhà máy X trên thị trường.
	\begin{enumerate}[a)]
		\item Tính xác suất để linh kiện điện tử đó được đóng dấu OTK.
		\item Dùng sơ đồ hình cây, hãy mô tả cách tính xác suất để linh kiện điện tử được chọn không được đóng dấu OTK.
	\end{enumerate}
	
	\loigiai{Gọi $A$ là biến cố: \lq \lq Linh kiện điện tử đó đạt tiêu chuẩn\rq \rq~. Gọi $B$ là biến cố: \lq \lq Linh kiện điện tử đó được đóng dấu OTK\rq \rq~.
		\begin{enumerate}[a)]
			\item Ta có $\mathrm{P}(B)=\mathrm{P}(A)\cdot \mathrm{P}(B \mid A)+\mathrm{P}(\overline{A}) \cdot \mathrm{P}(B \mid \overline{A})$.
			\begin{itemize}
				\item Tính $\mathrm{P}(A)$: Đây là xác suất để linh kiện đó đạt tiêu chuẩn. Vậy $\mathrm{P}(A)=0{,}8$.
				\item Tính $\mathrm{P}(\overline{A})$: $\mathrm{P}(\overline{A})=1-\mathrm{P}(A)=0{,}2$.
				\item Tính $\mathrm{P}(B\mid A)$: Đây là xác suất để linh kiện điện tử đó được đóng dấu OTK với điều kiện nó đạt tiêu chuẩn. Vậy $\mathrm{P}(B\mid A)=0{,}99$.
				\item Tính $\mathrm{P}(B\mid \overline{A})$: Đây là xác suất để linh kiện điện tử đó được đóng dấu OTK với điều kiện nó không đạt tiêu chuẩn. Vậy $\mathrm{P}(B\mid \overline{A})=1-0{,}95=0{,}05$.
			\end{itemize}
			Vậy $\mathrm{P}(B)=\mathrm{P}(A)\cdot \mathrm{P}(B \mid A)+\mathrm{P}(\overline{A}) \cdot \mathrm{P}(B \mid \overline{A})=0{,}8\cdot 0{,}99+0{,}2\cdot 0{,}05=0{,}802$. \\
			Vậy xác suất để  linh kiện điện tử đó được đóng dấu OTK là $0{,}802$.
			\item Ta có sơ đồ hình cây
			\begin{center}
				\begin{tikzpicture}[scale=0.8, font=\footnotesize, line join=round, line cap=round, >=stealth]
					\draw (0,0)--(-2,-2.5) node [midway, shift=(140:4mm)] {$0{,}8$} --(-3,-5) node [midway, shift=(160:4mm)] {$0{,}99$};
					\draw (0,0)--(2,-2.5) node [midway, shift=(40:4mm)] {$0{,}2$} --(3,-5) node [midway, shift=(20:4mm)] {$0{,}95$};
					\draw (-2,-2.5)--(-1,-5) node [midway, shift=(20:4mm)] {$0{,}01$};
					\draw (2,-2.5)--(1,-5) node [midway, shift=(160:4mm)] {$0{,}05$};
					\draw[draw=none,fill=yellow] (0,0) circle [radius=7pt] node[shift=(90:6mm)] {Gốc $O$};
					\draw[draw=none,fill=cyan] (-2,-2.5) circle [radius=7pt] node[shift=(140:6mm)] {$A$};
					\draw[draw=none,fill=green] (2,-2.5) circle [radius=7pt] node[shift=(40:6mm)] {$\overline{A}$};
					\draw[draw=none,fill=magenta] (1,-5) circle [radius=7pt] node[shift=(-90:6mm)] {$B$};
					\draw[draw=none,fill=magenta!40!black!30] (-1,-5) circle [radius=7pt] node[shift=(-90:6mm)] {$\overline{B}$};
					\draw[draw=none,fill=magenta] (-3,-5) circle [radius=7pt] node[shift=(-90:6mm)] {$B$};
					\draw[draw=none,fill=magenta!40!black!30] (3,-5) circle [radius=7pt] node[shift=(-90:6mm)] {$\overline{B}$};
				\end{tikzpicture}
			\end{center}
			Có hai nhánh cây đi từ $O$ tới $\overline{B}$ là $OA\overline{B}$ và $O\overline{A}\cdot\overline{B}$. Vậy xác suất để linh kiện điện tử được chọn không được đóng dấu OTK là $$\mathrm{P}(\overline{B})=0{,}8\cdot 0{,}01+0{,}2\cdot 0{,}95=0{,}198.$$
		\end{enumerate}
	}
\end{bt}

\begin{bt}%[2D6H2-2]
% nguon: Archive/2024/TaiLieuGiangDay/6. K12 - 4in 1 - XsCoDieuKien/data/CHUONG 6 - XAC SUAT/Bai2-Dang1-CongThucXSTP.tex #3284
[CKP, Mức độ 3]%[Trần Ngọc Thành]
	Trong một cuộc khảo sát tình trạng công việc trên $900$ người đã có bằng tốt nghiệp trung học phổ thông ở một địa phương cho cả nam lẫn nữ, người ta thu được số liệu thống kê trong bảng sau.
	\begin{center}
			\begin{tabular}{|l|c|c|}
			\hline
			\diagbox{Giới tính}{Tình trạng} & Có việc làm & Thất nghiệp\\\hline
			Nam & $460$ & $40$\\\hline
			Nữ & $140$ & $260$\\\hline
		\end{tabular}
	\end{center}
	Chọn ngẫu nhiên một người trong nhóm này. Gọi $A$ là biến cố \lq \lq Người được chọn là nữ\rq \rq , $B$ là biến cố \lq \lq Người được chọn có việc làm\rq \rq .
	\begin{enumerate}
		\item Vẽ lại sơ đồ hình cây sau đây và hoàn thành kết quả ở các ô \fbox{?}.
		\begin{center}
			\begin{tikzpicture}[line join = round, line cap = round, >=stealth, font=\footnotesize, scale=1]
				\begin{scope}[every node/.style={draw, rounded corners=5pt}]
					\node (A) at (0,0){Chọn một người};
					\def \gocA{30}
					\def \kcA{4}
					\node (B1) at ($(A)+(\gocA:\kcA)$){$A$};
					\node (B2) at ($(A)+(-\gocA:\kcA)$){$\overline{A}$};
					\def \gocB{15}
					\def \kcB{4}
					\node (B11) at ($(B1)+(\gocB:\kcB)$){$B$};
					\node (B12) at ($(B1)+(-\gocB:\kcB)$){$\overline{B}$};
					\node (B21) at ($(B2)+(\gocB:\kcB)$){$B$};
					\node (B22) at ($(B2)+(-\gocB:\kcB)$){$\overline{B}$};
				\end{scope}
				\begin{scope}[every node/.style={midway,sloped},every path/.style={->}]
					\draw (A)--(B1) node[above]{$\mathrm{P}(A)=$\fbox{?}};
					\draw (A)--(B2) node[below]{$\mathrm{P}(\overline{A})=$\fbox{?}};
					\draw (B1)--(B11) node[above]{$\mathrm{P}(B|A)=$\fbox{?}};
					\draw (B1)--(B12) node[below]{$\mathrm{P}(\overline{B}|A)=$\fbox{?}};
					\draw (B2)--(B21) node[above]{$\mathrm{P}(B|\overline{A})=$\fbox{?}};
					\draw (B2)--(B22) node[below]{$\mathrm{P}(\overline{B}|\overline{A})=$\fbox{?}};
				\end{scope}
				\def \kcC{1.7}
				\foreach \i/\j in {B11/AB,B12/{A\overline{B}},B21/{\overline{A}B},B22/{\overline{A}\,\overline{B}}}% Tạo nội dung lặp
				{
					\node at ($(\i)+(\kcC,0)$)[]{$\j$};
					\node at ($(\i)+({2*\kcC},0)$)[]{\fbox{?}};
				}
				\node (B) at ($(B11)+(\kcC,0.7)$){\textbf{Kết quả}};
				\node (C) at ($(B)+(\kcC,0)$){\textbf{Xác suất}};
			\end{tikzpicture}\\
			$A \colon$ nữ; $\overline{A} \colon$ nam; $B \colon$ có việc; $\overline{B} \colon$ thất nghiệp.
		\end{center}
		\item Tính xác suất để chọn được một người có việc làm.
	\end{enumerate}
	\loigiai{
		\begin{enumerate}
			\item Theo đề bài xác suất để chọn được một người nữ là $\mathrm{P}(A)=\dfrac{4}{9}$, suy ra $\mathrm{P}(\overline{A})=\dfrac{5}{9}$.\\
			Xác suất chọn được người có việc làm nếu người đó là nữ $\mathrm{P}(B|A)=\dfrac{140}{400}=\dfrac{7}{20}$. Suy ra $\mathrm{P}(\overline{B}|A)=\dfrac{13}{20}$.\\
			Xác suất chọn được người có việc làm nếu người đó không là nữ $\mathrm{P}(B|\overline{A})=\dfrac{460}{500}=\dfrac{23}{25}$.\\
			Suy ra $\mathrm{P}(\overline{B}|\overline{A})=\dfrac{2}{25}.$\\
			\begin{center}
				\begin{tikzpicture}[line join = round, line cap = round, >=stealth, font=\footnotesize, scale=1]
					\begin{scope}[every node/.style={draw, rounded corners=5pt}]
						\node (A) at (0,0){Chọn một người};
						\def \gocA{30}
						\def \kcA{4}
						\node (B1) at ($(A)+(\gocA:\kcA)$){$A$};
						\node (B2) at ($(A)+(-\gocA:\kcA)$){$\overline{A}$};
						\def \gocB{15}
						\def \kcB{4}
						\node (B11) at ($(B1)+(\gocB:\kcB)$){$B$};
						\node (B12) at ($(B1)+(-\gocB:\kcB)$){$\overline{B}$};
						\node (B21) at ($(B2)+(\gocB:\kcB)$){$B$};
						\node (B22) at ($(B2)+(-\gocB:\kcB)$){$\overline{B}$};
					\end{scope}
					\begin{scope}[every node/.style={midway,sloped},every path/.style={->}]
						\draw (A)--(B1) node[above]{$\mathrm{P}(A)=$\fbox{$\dfrac{4}{9}$}};
						\draw (A)--(B2) node[below]{$\mathrm{P}(\overline{A})=$\fbox{$\dfrac{5}{9}$}};
						\draw (B1)--(B11) node[above]{$\mathrm{P}(B|A)=$\fbox{$\dfrac{7}{20}$}};
						\draw (B1)--(B12) node[below]{$\mathrm{P}(\overline{B}|A)=$\fbox{$\dfrac{13}{20}$}};
						\draw (B2)--(B21) node[above]{$\mathrm{P}(B|\overline{A})=$\fbox{$\dfrac{23}{25}$}};
						\draw (B2)--(B22) node[below]{$\mathrm{P}(\overline{B}|\overline{A})=$\fbox{$\dfrac{2}{25}$}};
					\end{scope}
					\def \kcC{1.7}
					\foreach \i/\j/\k in {B11/AB/{\dfrac{7}{45}},B12/{A\overline{B}}/{\dfrac{13}{45}},B21/{\overline{A}B}/{\dfrac{23}{45}},B22/{\overline{A}\,\overline{B}}/{\dfrac{2}{45}}}% Tạo nội dung lặp
					{
						\node at ($(\i)+(\kcC,0)$)[]{$\j$};
						\node at ($(\i)+({2*\kcC},0)$)[]{$\k$};
					}
					\node (B) at ($(B11)+(\kcC,0.7)$){\textbf{Kết quả}};
					\node (C) at ($(B)+(\kcC,0)$){\textbf{Xác suất}};
				\end{tikzpicture}
			\end{center} 
			\item Xác suất để chọn được một người có việc làm $$\mathrm{P}(B)=\mathrm{P}(A)\mathrm{P}(B|A)+\mathrm{P}(\overline{A})\mathrm{P}(B|\overline{A})=\dfrac{4}{9}\cdot \dfrac{7}{20}+\dfrac{5}{9}\cdot \dfrac{23}{25}=\dfrac{2}{3}.$$
		\end{enumerate}
	}
\end{bt}

\begin{bt}%[2D6V2-2]
% nguon: Archive/2024/TaiLieuGiangDay/6. K12 - 4in 1 - XsCoDieuKien/data/CHUONG 6 - XAC SUAT/Bai2-Dang1-CongThucXSTP.tex #3285
[CTST, Mức độ 3]%[Trần Ngọc Thành]
	Hộp thứ nhất có $1$ viên bi xanh và $5$ viên bi đỏ. Hộp thứ hai có $3$ viên bi xanh và $5$ viên bi đỏ. Các viên bi có cùng kích thước và khối lượng. Lấy ra ngẫu nhiên đồng thời $2$ viên bi từ hộp thứ nhất chuyển sang hộp thứ hai. Sau đó lại lấy ra ngẫu nhiên $2$ viên bi từ hộp thứ hai.
  Tính xác suất để hai viên bi lấy ra từ hộp thứ hai là bi đỏ.
	\loigiai{
	Xét các biến cố
			\begin{itemize}
				\item $A$: \lq\lq Hai viên bi lấy ra từ hộp thứ hai là bi đỏ\rq\rq.
				\item $B_1$: \lq\lq Hai viên bi lấy ra từ hộp thứ nhất có cả màu xanh và màu đỏ\rq\rq. 
				\item $B_2$: \lq\lq Hai viên bi lấy ra từ hộp thứ nhất có màu đỏ\rq\rq.
			\end{itemize}
			Ta có 
			\[\mathrm{P}(B_1) = \dfrac{\mathrm{C}^1_{5}}{\mathrm{C}^2_{6}}=\dfrac{1}{3};\quad 
			\mathrm{P}(B_2) = \dfrac{\mathrm{C}^2_5}{\mathrm{C}^2_{6}}=\dfrac{2}{3};\quad
			\mathrm{P}(A\mid B_1) = \dfrac{\mathrm{C}^2_6}{\mathrm{C}^2_{10}}=\dfrac{1}{3};\quad
			\mathrm{P}(A\mid B_2) = \dfrac{\mathrm{C}^2_7}{\mathrm{C}^2_{10}}=\dfrac{7}{15}.
			\]
			Áp dụng công thức xác suất toàn phần, ta có
			\allowdisplaybreaks
			\begin{eqnarray*}
				\mathrm{P}(A) 
				&=& \mathrm{P}(A|B_1)\cdot \mathrm{P}(B_1) + \mathrm{P}(A|B_2)\cdot\mathrm{P}(B_2)\\ 
				&=& \dfrac{1}{3}\cdot \dfrac{1}{3} + \dfrac{7}{15}\cdot \dfrac{2}{3}\\
				&=& \dfrac{19}{45}.
			\end{eqnarray*}
		
}\end{bt}

\begin{bt}%[2D6H2-2]
% nguon: Archive/2024/TaiLieuGiangDay/6. K12 - 4in 1 - XsCoDieuKien/data/CHUONG 6 - XAC SUAT/Bai2-Dang1-CongThucXSTP.tex #3286
[CD, Mức độ 4]%[Trần Ngọc Thành]
	Có hai chiếc hộp, hộp I có $5$ viên bi màu trắng và $5$ viên bi màu đen; hộp II có $6$ viên bi màu trắng và $4$ viên bi màu đen. Các viên bi có cùng kích thước và khối lượng. Lấy ngẫu nhiên đồng thời hai viên bi từ hộp I bỏ sang hộp II. Sau đó lấy ngẫu nhiên một viên bi từ hộp II.
 Tính xác suất để viên bi được lấy ra là viên bi màu trắng.
	
	\loigiai{
 Xét các biến cố
			\begin{itemize}
				\item[] $A$: \lq\lq Viên bi lấy ra là viên màu trắng\rq\rq.
				\item[] $B_1$: \lq\lq 2 viên bi lấy ra từ hộp I có màu trắng\rq\rq. 
				\item[] $B_2$: \lq\lq 2 viên bi lấy ra từ hộp I có màu đen\rq\rq.
				\item[] $B_3$: \lq\lq Cả 2 viên bi lấy ra từ hộp I có cả hai màu đen trắng\rq\rq.
			\end{itemize}
			Ta có 
			\[\mathrm{P}(B_1) = \dfrac{\mathrm{C}^2_{5}}{\mathrm{C}^2_{10}}=\dfrac{2}{9};\quad 
			\mathrm{P}(B_2) = \dfrac{\mathrm{C}^2_5}{\mathrm{C}^2_{10}}=\dfrac{2}{9};\quad
			\mathrm{P}(B_3) = \dfrac{\mathrm{C}^1_5\cdot\mathrm{C}^1_5}{\mathrm{C}^2_{10}}=\dfrac{5}{9}.
			\]
			Áp dụng công thức xác suất toàn phần, ta có
			\allowdisplaybreaks
			\begin{eqnarray*}
				\mathrm{P}(A) 
				&=& \mathrm{P}(A|B_1)\cdot \mathrm{P}(B_1) + \mathrm{P}(A|B_2)\cdot\mathrm{P}(B_2) + \mathrm{A|B_3}\cdot\mathrm{P}(B_3)\\ 
				&=& \dfrac{\mathrm{C}^2_8}{\mathrm{C}^2_{12}}\cdot \dfrac{2}{9} + \dfrac{\mathrm{C}^2_{6}}{\mathrm{C}^2_{12}}\cdot \dfrac{2}{9} + \dfrac{\mathrm{C}^2_7}{\mathrm{C}^2_{12}}\cdot \dfrac{5}{9}\\
				&=& \dfrac{14}{33}\cdot \dfrac{2}{9} + \dfrac{5}{22}\cdot \dfrac{2}{9} + \dfrac{7}{22}\cdot \dfrac{5}{9}\\
				&=& \dfrac{191}{594}.
			\end{eqnarray*}

	}
\end{bt}

\begin{bt}%[2D6V2-2]
% nguon: Archive/SP-4Phan-Dot3/Data/CD_2-OTHK2-De01.tex #4421
%[Dự án đè thi 4 phần đợt 3]%[Trần Xuân Hòa]
	Một đội văn nghệ gồm $n$ bạn nam và $7$ bạn nữ. Chọn ra ngẫu nhiên $2$ bạn để biểu diễn một tiết mục. Gọi $A$ là biến cố: \lq\lq Có ít nhất một bạn nữ trong $2$ bạn được chọn\rq\rq, $B$ là biến cố \lq\lq Hai bạn được chọn có cùng giới tính\rq\rq. Biết xác suất của biến cố $A$ với điều kiện $B$ bằng $\dfrac{21}{31}$, tìm số bạn nam $n$ của đội văn nghệ.
	\loigiai{
		Ta có $\mathrm{P}(B)=\dfrac{\mathrm{C}_n^2+\mathrm{C}_7^2}{\mathrm{C}_{n+7}^2}=\dfrac{\mathrm{C}_n^2+21}{\mathrm{C}_{n+7}^2}$. \\
			Vì $A B$ là biến cố \lq\lq Hai bạn được chọn đều là nữ\rq\rq \,nên $\mathrm{P}(A B)=\dfrac{\mathrm{C}_7^2}{\mathrm{C}_{n+7}^2}=\dfrac{21}{\mathrm{C}_{n+7}^2}$. \\
			Do đó	$\mathrm{P}(A \mid B)=\dfrac{\mathrm{P}(A B)}{\mathrm{P}(B)}=\dfrac{\dfrac{21}{\mathrm{C}_{n+7}^2}}{\dfrac{\mathrm{C}_n^2+21}{\mathrm{C}_{n+7}^2}}=\dfrac{21}{\mathrm{C}_n^2+21} $. \\
			Theo giả thiết ta có $\mathrm{P}(A \mid B)=\dfrac{21}{31}$ hay $\dfrac{21}{\mathrm{C}_n^2+21}=\dfrac{21}{31}$. \\
			Do đó $\mathrm{C}_n^2=10$ hay $n(n-1)=20$. \\
			Vì $n$ là số nguyên dương nên ta tìm được giá trị của $n$ bằng $5$.
		}
\end{bt}

\begin{bt}%[2D6H2-1]
% nguon: Archive/SP-4Phan-Dot3/Data/KNTT_2-OTHK2_De02.tex #4462

	Cho hai biến cố $A$ và $B$ với $0<\mathrm{P}\left( A \right)<1$.
	\begin{enumerate}
		\item Viết công thức xác suất toàn phần tính $\mathrm{P}\left( B \right)$.
		\item Với 
		$\mathrm{P}\left( A \right)=0{,}1$; 
		$\mathrm{P}\left( \overline{A} \right)=0{,}9$; 
		$\mathrm{P}\left( B|A \right)=0{,}3$; 
		$\mathrm{P}\left( B|\overline{A} \right)=0{,}6$. 
		Tính $\mathrm{P}\left( B \right)$?
	\end{enumerate}  
	\loigiai{
		\begin{enumerate}
			\item Viết công thức xác suất toàn phần tính $\mathrm{P}\left( B \right)$.\\
			Ta có: $\mathrm{P}\left( B \right)=
			\mathrm{P}\left( A \right)\mathrm{P}\left( B|A \right)+
			\mathrm{P}\left( \overline{A} \right)\mathrm{P}\left( B|\overline{A} \right)$.
			\item Với 
			$\mathrm{P}\left( A \right)=0{,}1$; 
			$\mathrm{P}\left( \overline{A} \right)=0{,}9$; 
			$\mathrm{P}\left( B|A \right)=0{,}3$; 
			$\mathrm{P}\left( B|\overline{A} \right)=0{,}6$. 
			Tính $\mathrm{P}\left( B \right)$?\\
			Áp dụng công thức xác suất toàn phần:\\
			$\mathrm{P}\left( B \right)=
			\mathrm{P}\left( A \right)\mathrm{P}\left( B|A \right)+
			\mathrm{P}\left( \overline{A} \right)\mathrm{P}\left( B|\overline{A} \right)=
			0{,}1\cdot 0{,}3+0{,}9\cdot 0{,}6=0{,}57$.			
		\end{enumerate}			
	}	
\end{bt}

\begin{bt}%[2D6H1-2]
% nguon: Archive/SP-4Phan-Dot3/Data/CD_2-OTHK2_De04.tex #4499
%[Nguyễn Tuấn, sáng tác đề GHK2, HK2 lần 6]
	Một sản phẩm xuất khẩu sang Mỹ phải qua $2$ lần kiểm tra, nếu cả $2$ lần đều đạt thì sản phẩm đó mới đủ tiêu chuẩn xuất khẩu. Biết rằng bình quân $98\%$ sản phẩm làm ra qua được lần kiểm tra thứ nhất và $95\%$ số sản phẩm qua lần kiểm tra đầu tiên sẽ tiếp tục qua được lần kiểm tra thứ hai. Xác suất để sản phẩm đó đủ tiêu chuẩn xuất khẩu là $\dfrac{a}{b}$ (với $a$, $b$ là các số nguyên dương và $\dfrac{a}{b}$ là phân số tối giản). Giá trị của $b-a$ bằng bao nhiêu?
	\loigiai{
		Gọi $A$ là biến cố \lq\lq qua được lần kiểm tra đầu tiên\rq\rq, suy ra $\mathrm{P}(A)=0{,}98$.\\
		Gọi $B$ là biến cố \lq\lq qua được lần kiểm tra thứ hai\rq\rq, suy ra $\mathrm{P}(B\mid A)=0{,}95$.\\
		Sản phẩm đủ tiêu chuẩn xuất khẩu phải qua được cả $2$ lần kiểm tra nên ta cần tính $\mathrm{P}(AB)$.\\
		Ta có $\mathrm{P}(B\mid A)=\dfrac{\mathrm{P}(AB)}{\mathrm{P}(A)} \Rightarrow \mathrm{P}(AB)=\mathrm{P}(A)\cdot \mathrm{P}(B\mid A)=0{,}98\cdot 0{,}95=0{,}931=\dfrac{931}{1000}$.\\
		Khi đó $a=931$, $b=1000\Rightarrow b-a=69$.
	}
\end{bt}

\begin{bt}%[2D6H1-2]
% nguon: Archive/SP-4Phan-Dot3/Data/CTST_2-OTHK2-De05.tex #4547
%Câu 17
Một nhà máy có hai phân xưởng I và II. Phân xưởng I sản xuất $40\rm{\%}$ số sản phẩm và phân xưởng $\rm{II}$ sản xuất $60\rm{\%}$ số sản phẩm. Tỉ lệ sản phẩm bị lỗi của phần xưởng $\rm{I}$ là $3\rm{\%}$ và của phân xưởng II là $\rm{2\%}$. Kiểm tra ngẫu nhiên 1 sản phẩm của nhà máy và tính xác suất để sản phẩm đó bị lỗi là $\dfrac{t}{100}$. Khi đó $t$ bằng bao nhiêu?
\loigiai{
Gọi $A$ là biến cố \lq\lq Sản phẩm bị lỗi\rq\rq\, và $B$ là biến cố \lq\lq Sản phẩm lấy ra do phân xường I sản xuất\rq\rq.\\
Do phân xưởng I sản xuất $40\%$ số sản phẩm và phân xưởng II sản xuất $60\%$ số sản phẩm nên $\mathrm{P}(B)=0{,}4$, $\mathrm{P}\left(\overline{B}\right)=1-0{,}4=0{,}6.$\\
Do tỉ lệ sản phẩm bị lỗi của phân xưởng I là $2\%$ và của phân xưởng II là $1\%$ nên
\[\mathrm{P}\left(A| B\right)=0{,}03; \, \mathrm{P}\left(A|\overline{B}\right)=0{,}02.\]
Xác suất để sản phẩm lấy ra bị lỗi là
\allowdisplaybreaks
\begin{eqnarray*}
&&\mathrm{P}(A)=\mathrm{P}(B)\mathrm{P}\left(A| B\right)+\mathrm{P}\left(\overline{B}\right)\mathrm{P}\left(A|\overline{B}\right)\\
&=&0{,}4\cdot0{,}03+0{,}6\cdot0{,}02=0{,}024=\dfrac{2{,}4}{100}.
\end{eqnarray*}
Vậy $t=2{,}4$.
}
\end{bt}

\begin{bt}%[2D6H1-2]
% nguon: Archive/SP-4Phan-Dot3/Data/CTST_2-OTHK2-De04.tex #4554
 
	Cho hai biến cố $A$, $B$ có $P(A)=0{,}8$; $P(B)=0{,}5$; $P(A \cap B)=0{,}2$. Tính xác suất biến cố $B$ không xảy ra với điều kiện biến cố $A$ xảy ra.
	\loigiai{
		Xác suất biến cố $B$ không xảy ra với điều kiện biến cố $A$ xảy ra là
		$$P(\overline{B} |A)=\dfrac{P(A\overline{B})}{P(A)}=\dfrac{P(A)-P(A\cap B)}{P(A)}=\dfrac{3}{4}.$$
	}
\end{bt}

\begin{bt}%[2D6C1-4]
% nguon: Archive/SP-4Phan-Dot3/Data/CTST_2-OTHK2-De04.tex #4557
 
	Kết quả một cuộc khảo sát các vụ tai nạn giao thông ô tô về mối quan hệ giữa việc thắt dây an toàn của người lái xe khi xảy ra tai nạn giao thông và nguy cơ tử vong của người lái xe khi xảy ra tai nạn giao thông cho thấy:
	\begin{itemize}
		\item Tỉ lệ người lái xe tử vong khi xảy ra tai nạn giao thông là $0{,4}\%$.
		\item Tỉ lệ người lái xe không thắt dây an toàn giao thông khi xảy ra tai nạn giao thông là $28 \%$.
		\item Tỉ lệ người lái xe tử vong khi xảy ra tai nạn giao thông trong trường hợp không thắt dây an toàn là $0{,}3 \%$.
	\end{itemize}
	Hỏi theo kết quả khảo sát trên, việc thắt dây an toàn của người lái xe ô tô sẽ làm giảm khả năng tử vong là bao nhiêu lần?
	\loigiai{
		Chọn ngẫu nhiên một một vụ tai nạn giao thông của cuộc khảo sát trên. Xét các biến cố:\\
		$A$: "Người lái xe đó tử vong khi xảy ra tai nạn giao thông."\\
		$B$: "Người lái xe đó không thắt dây an toàn khi xảy ra tai nạn giao thông."\\
		Ta có $P(A)= 0{,4}\%$; $P(B)= 28\%$; $P(A\cap B)=0{,}3\%$.\\
		Xác suất người lái xe đó tử vong khi xảy ra tai nạn giao thông trong trường hợp không thắt dây an toàn là
		$$P(A|B)=\dfrac{P(A\cap B)}{P(B)}=\dfrac{3}{280}.$$
		Xác suất người lái xe đó có thắt dây an toàn giao thông là $P(\overline{B})=72\%$.\\
		Xác suất người lái xe đó tử vong khi xảy ra tai nạn giao thông trong trường hợp có thắt dây an toàn là
		$$P(A|\overline{B})=\dfrac{P(A\cap \overline{B})}{P(\overline{B})}=\dfrac{P(A)-P(A\cap B)}{P(\overline{B}}=\dfrac{1}{720}.$$
		Ta có
		$$\dfrac{P(A|B)}{P(A|\overline{B})}=\dfrac{54}{7}\approx7{,}7.$$
		Vậy theo khảo sát trên, việc thắt dây an toàn của người lái xe ô tô sẽ làm giảm khả năng tử vong khoảng $7{,}7$ lần.
	}
\end{bt}

\begin{bt}%[2D6H1-2]
% nguon: Archive/SP-4Phan-Dot3/Data/CTST_2-OTHK2-De06.tex #4562

Gieo một con xúc xắc cân đối và đồng chất. Tính xác suất số chấm trên con xúc xắc không nhỏ hơn $4$, biết rằng con xúc xắc xuất hiện mặt lẻ.
\loigiai
{
Gọi $A$ là biến cố: \lq\lq Số chấm trên xúc xắc không nhỏ hơn $4$\rq\rq.\\
$B$ là biến cố: \lq\lq Xúc xắc xuất hiện mặt lẻ\rq\rq.\\
Ta cần tính $\mathrm{P}(A| B)$.\\
Kết quả thuận lợi của biến cố $A=\{4;5;6\}$.\\
Kết quả thuận lợi của biến cố $B=\{1;3;5\}$.\\
Suy ra $A\cap B=\{5\}$.\\
Suy ra $\heva{
&\mathrm{P}(A\cap B)=\dfrac{1}{6}\\
&\mathrm{P}(B)=\dfrac{3}{6}.
}$\\
Vậy $\mathrm{P}(A\mid B)=\dfrac{\dfrac{3}{6}}{\dfrac{1}{6}}=\dfrac{1}{3}$.
}
\end{bt}

\begin{bt}%[2D6H2-3]
% nguon: Archive/SP-4Phan-Dot3/Data/CTST_2-OTHK2-De06.tex #4564

	Trong $1$ đám đông, số người nam bằng số người nữ. Xác suất mắc cận thị của nam là $0{,}4$ và nữ là $0{,}6$. Chọn ngẫu nhiên $1$ người. Xác suất chọn được nam không cận thị (làm tròn đến kết quả hàng phần trăm).
	\loigiai{ 
		Goi $H_1$: \lq\lq Biến cố chọn nam\rq\rq; $H_2$ là biến cố: \lq\lq Chọn được nữ\rq\rq.\\
		$C$: \lq\lq Biến cố người được chọn là không cận thị\rq\rq.\\
		Ta có $\mathrm{P}\left(H_1\right)=\mathrm{P}\left(H_2\right)=\dfrac{1}{2}$\\
		$\mathrm{P}(C)=\mathrm{P}\left(H_1\right)\mathrm{P}\left(C\mid H_1\right)+\mathrm{P}\left(H_2\right)\mathrm{P}\left(C\mid H_2\right)=\dfrac{1}{2}\cdot0{,}6+\dfrac{1}{2}\cdot0{,}4=0{,}5$.\\
		Xác suất để chọn ngẫu nhiên ra $1$ người mà người đó là nam không cận thị là\\ (Áp dụng định lý Bayes)\\
		\[\mathrm{P}\left(H_1\mid C\right)=\dfrac{\mathrm{P}\left(H_1\right)\mathrm{P}\left(C\mid H_1\right)}{\mathrm{P}(C)}=\dfrac{\dfrac{1}{2}\cdot0{,}6}{0{,}5}=0{,}6.\]
		}
\end{bt}

\begin{bt}%[2D6V1-2]
% nguon: Archive/SP-4Phan-Dot3/Data/CTST_2-OTHK2-De03.tex #4573

	Trong hộp có $3$ viên bi màu trắng và $7$ viên bi màu đỏ. Lấy lần lượt mỗi lần một viên theo cách lấy không trả lại.\\
	Gọi $A$ là biến cố \lq\lq viên bi lấy lần thứ nhất là màu đỏ\rq\rq.\\
	Gọi $B$ là biến cố \lq\lq viên bi lấy lần thứ hai là màu đỏ\rq\rq.\\
	Tính $\mathrm{P}(A\overline{B})$.
	\loigiai{
		Dễ thấy rằng $\mathrm{P}(A)=0{,}7$.\\
		Ta có $\mathrm{P}(B\mid A)=\dfrac{6}{9}=\dfrac{2}{3}$ và $\mathrm{P}(B\mid \overline{A})=\dfrac{7}{9}$.\\
		Do $\mathrm{P}(A)=0{,}7$ nên $\mathrm{P}(\overline{A})=0{,}3$.\\
		Suy ra, $\mathrm{P}(B\overline{A})=\mathrm{P}(\overline{A})\cdot \mathrm{P}(B\mid \overline{A})=\dfrac{7}{30}$.\\
		$\mathrm{P}(B)=\mathrm{P}(A)\cdot \mathrm{P}(B|A)+\mathrm{P}({\overline{A}})\cdot \mathrm{P}(B|\overline{A})=0{,}7\cdot\dfrac{2}{3}+\dfrac{7}{9}\cdot 0{,}3=0{,}7$.\\
		$\mathrm{P}(\overline{B}A)=\dfrac{\mathrm{P}(\overline{B}\mid A)}{ \mathrm{P}(A)}=\dfrac{1-\mathrm{P}(B\mid A)}{\mathrm{P}(A)}=\dfrac{10}{21}$.
	} 
\end{bt}

\begin{bt}%[2D6H1-2]
% nguon: Archive/SP-4Phan-Dot3/Data/CTST_2-OTHK2-De03.tex #4574

	Một công ty bảo hiểm nhận thấy có $48\%$ số người mua bảo hiểm ô tô là phụ nữ và có $36\%$ số người mua bảo hiểm ô tô là phụ nữ trên $45$ tuổi. Biết một người mua bảo hiểm ô tô là phụ nữ, tính xác suất người đó trên $45$ tuổi.
	\loigiai{
		Gọi $A$ là biến cố \lq\lq Người mua bảo hiểm ô tô là phụ nữ\rq\rq, $B$ là biến cố \lq\lq Người mua bảo hiểm ô tô trên $45$ tuổi\rq\rq. Ta cần tính $\mathrm{P}(B \mid A)$.\\
		Do có $48 \%$ người mua bảo hiểm ô tô là phụ nữ nên $\mathrm{P}(A)=0{,}48$.\\
		Do có $36 \%$ số người mua bảo hiểm ô tô là phụ nữ trên $45$ tuổi nên $\mathrm{P}(A B)=0{,}36$.\\
		Vậy $\mathrm{P}(B \mid A)=\dfrac{\mathrm{P}(A B)}{\mathrm{P}(A)}=\dfrac{0{,}36}{0{,}48}=0{,}75$.  
	}
\end{bt}

\begin{bt}%[2D6V2-2]
% nguon: Archive/SP-4Phan-Dot3/Data/CTST_2-OTHK2-De02.tex #4582

    Có hai hộp đựng bóng. Hộp I có $4$ quả bóng màu xanh và $8$ quả bóng màu đỏ. Hộp II có $6$ quả bóng màu xanh và $4$ quả bóng màu đỏ. Trước tiên, từ hộp II lấy ra ngẫu nhiên $1$ quả bóng rồi cho vào hộp I. Sau đó, từ hộp I lấy ra ngẫu nhiên $1$ quả bóng. Tính xác suất để quả bóng được lấy ra từ hộp I là quả bóng màu đỏ (kết quả làm tròn đến hàng phần trăm).
    \loigiai{
        Gọi
        \begin{itemize}
            \item Biến cố $A\colon$ \lq\lq Quả bóng được lấy ra từ hộp II để cho vào hộp I là quả bóng màu đỏ\rq\rq.
            \item Biến cố $B\colon$ \lq\lq Quả bóng được lấy ra từ hộp I là quả bóng màu đỏ\rq\rq.
        \end{itemize}
        Ta cần tính xác suất của biến cố $B$ là $\mathrm{P}(B)$.\\
        Sử dụng công thức xác suất toàn phần
        \[\mathrm{P}(B) = \mathrm{P}(A) \cdot \mathrm{P}(B\mid A) + \mathrm{P}\left(\overline{A}\right) \cdot \mathrm{P}\left(B\mid\overline{A}\right).\]
        Tính các xác suất thành phần
        \begin{itemize}
            \item Xác suất $\mathrm{P}(A)$: Đây là xác suất để lấy ra ngẫu nhiên $1$ quả bóng đỏ từ hộp II rồi cho vào hộp I. Hộp II có tổng cộng $10$ quả bóng, trong đó có $4$ quả bóng đỏ. Do đó
            \[\mathrm{P}(A) = \dfrac{4}{10} = 0{,}4.\]
            \item Xác suất $\mathrm{P}\left(\overline{A}\right)$: Đây là xác suất lấy ra $1$ quả bóng xanh từ hộp II rồi cho vào hộp I. Khi đó
            \[\mathrm{P}\left(\overline{A}\right) = 1 - \mathrm{P}(A) = 1 - 0{,}4 = 0{,}6.\]
            \item Xác suất $\mathrm{P}(B\mid A)$: Đây là xác suất để lấy ra $1$ quả bóng đỏ từ hộp I, với điều kiện đã chọn $1$ quả bóng đỏ từ hộp II và cho vào hộp I. Khi đó, hộp I có $4$ quả bóng xanh và $9$ quả bóng đỏ. Xác suất để lấy ra một quả bóng đỏ là
            \[\mathrm{P}(B\mid A) = \dfrac{9}{13}.\]
            \item Xác suất $\mathrm{P}\left(B|\overline{A}\right)$: Đây là xác suất để lấy ra $1$ quả bóng đỏ từ hộp I, với điều kiện đã chọn $1$ quả bóng xanh từ hộp II và cho vào hộp I. Khi đó, hộp I có $5$ quả bóng xanh và $8$ quả bóng đỏ. Xác suất để lấy ra một quả bóng đỏ là
            \[\mathrm{P}\left(B\mid\overline{A}\right) = \dfrac{8}{13}.\]
        \end{itemize}
        Vậy
        \begin{eqnarray*}
        	\mathrm{P}(B) &=& \mathrm{P}(A) \cdot \mathrm{P}(B\mid A) + \mathrm{P}\left(\overline{A}\right) \cdot \mathrm{P}\left(B\mid\overline{A}\right)\\
        	&=& 0{,}4 \cdot \dfrac{9}{13} + 0{,}6 \cdot \dfrac{8}{13} = \dfrac{36 + 48}{130} = \dfrac{84}{130} \approx 0{,}65.
        \end{eqnarray*}
        Xác suất để quả bóng được lấy ra từ hộp I là quả bóng màu đỏ là $0{,}65$ (hay $65\%$).
    }
\end{bt}

\begin{bt}%[2D6H2-3]
% nguon: Archive/SP-4Phan-Dot3/Data/CTST_2-OTHK2-De02.tex #4583
 
    Một thống kê cho thấy tỉ lệ dân số mắc bệnh nhiễm trùng $X$ là $0{,}8\%$. Ông $M$ đi xét nghiệm bệnh nhiễm trùng $X$ và nhận được kết quả là dương tính. Biết rằng, nếu mắc bệnh nhiễm trùng $X$ thì với xác suất $0{,}92$ xét nghiệm cho kết quả dương tính; nếu không bị bệnh nhiễm trùng $X$ thì với xác suất $0{,}96$ xét nghiệm cho kết quả âm tính. Hỏi sau khi xét nghiệm cho kết quả dương tính, xác suất ông $M$ mắc bệnh nhiễm trùng $X$ là bao nhiêu?
    \loigiai{
        Gọi $A$ là biến cố: \lq\lq Ông $M$ bị bệnh nhiễm trùng X\rq\rq; \\
        $B$ là biến cố: \lq\lq Xét nghiệm cho kết quả dương tính\rq\rq.
        Ta có
        \begin{itemize}
            \item  $\mathrm{P}(A) = 0{,}008$ (tỉ lệ dân số mắc bệnh nhiễm trùng $X$).
            \item $\mathrm{P}(B\mid A) = 0{,}92$ (xác suất xét nghiệm dương tính khi mắc bệnh).
            \item $\mathrm{P}\left(\overline{B}\mid \overline{A}\right) = 0{,}96$ (xác suất xét nghiệm âm tính khi không mắc bệnh).
            \item $\mathrm{P}\left(\overline{A}\right) = 1 - \mathrm{P}(A) = 1 - 0{,}008 = 0{,}992$.
        \end{itemize}
            Sau khi xét nghiệm cho kết quả dương tính, xác suất ông $M$ mắc bệnh nhiễm trùng $X$ là $\mathrm{P}(A\mid B)$.\\
            Ta cần tính $\mathrm{P}(A\mid B)$ theo công thức Bayes
            \[\mathrm{P}(A\mid B) = \dfrac{\mathrm{P}(A) \cdot \mathrm{P}(B\mid A)}{\mathrm{P}(A) \cdot \mathrm{P}(B\mid A) + \mathrm{P}\left(\overline{A}\right) \cdot \mathrm{P}\left(B\mid \overline{A}\right)}.	\]
            Biết rằng $\mathrm{P}\left(B\mid \overline{A}\right) = 1 - \mathrm{P}\left(\overline{B}\mid\overline{A}\right) = 1 - 0{,}96 = 0{,}04$.
            Thay số vào, ta có
            \[\mathrm{P}(A\mid B) = \dfrac{0{,}008 \cdot 0{,}92}{0{,}008 \cdot 0{,}92 + 0{,}992 \cdot 0{,}04} = \dfrac{0{,}00736}{0{,}00736 + 0{,}03968} \approx 0{,}156.\]
    }
\end{bt}

\begin{bt}%[2D6V1-2]
% nguon: Archive/SP-4Phan-Dot3/Data/CTST_2-OTHK2-De01.tex #4592

	Người ta điều tra thấy ở một địa phương X có $2\%$ tài xế
	sử dụng điện thoại di động khi lái xe. Trong các vụ tai nạn ở
	địa phương đó, người ta nhận thấy có $10\%$ là do tài xế có sử dụng
	điện thoại khi lái xe gây ra. Hỏi việc sử dụng điện thoại
	di động khi lái xe làm tăng xác suất gây tai nạn lên bao nhiêu lần?
	\loigiai{Gọi $ A $ là biến cố: \lq\lq Tài xế sử dụng điện thoại di động khi lái xe\rq\rq\, và $ B $ là biến cố: \lq\lq Địa phương có tai nạn\rq\rq.\\
		Do có $ 2\% $ tài xế sử dụng điện thoại di động khi lái xe nên $ \mathrm{P}(A)=0{,}02 $.\\
		Do trong các vụ tai nạn ở địa phương, có $ 10\% $ là do tài xế có sử dụng điện thoại khi lái xe nên $ \mathrm{P}(A|B)=0,1 $.
		Xác suất gây tai nạn do tài xế sử dụng điện thoại di động khi lái xe là
		$$ \mathrm{P}(B|A)= \dfrac{\mathrm{P}(B)\mathrm{P}(A|B)}{\mathrm{P}(A)}.$$
		Do đó $ \dfrac{\mathrm{P}(B|A)}{\mathrm{P}(B)}=\dfrac{\mathrm{P}(A|B)}{\mathrm{P}(A)}=\dfrac{0{,}1}{0{,}02}= 5$.\\ Suy ra $ \mathrm{P}(B|A)=5\mathrm{P}(B) $.\\
		Vậy việc sử dụng điện thoại di động khi lái xe làm tăng xác suất gây tai nạn lên $ 5 $ lần.}	
\end{bt}
